% Options for packages loaded elsewhere
\PassOptionsToPackage{unicode}{hyperref}
\PassOptionsToPackage{hyphens}{url}
\PassOptionsToPackage{dvipsnames,svgnames,x11names}{xcolor}
%
\documentclass[
  11pt,
]{article}
\usepackage{amsmath,amssymb}
\usepackage{iftex}
\ifPDFTeX
  \usepackage[T1]{fontenc}
  \usepackage[utf8]{inputenc}
  \usepackage{textcomp} % provide euro and other symbols
\else % if luatex or xetex
  \usepackage{unicode-math} % this also loads fontspec
  \defaultfontfeatures{Scale=MatchLowercase}
  \defaultfontfeatures[\rmfamily]{Ligatures=TeX,Scale=1}
\fi
\usepackage{lmodern}
\ifPDFTeX\else
  % xetex/luatex font selection
  \setmainfont[]{Liberation Serif}
  \setsansfont[]{Liberation Sans}
  \setmonofont[]{DejaVu Sans Mono}
\fi
% Use upquote if available, for straight quotes in verbatim environments
\IfFileExists{upquote.sty}{\usepackage{upquote}}{}
\IfFileExists{microtype.sty}{% use microtype if available
  \usepackage[]{microtype}
  \UseMicrotypeSet[protrusion]{basicmath} % disable protrusion for tt fonts
}{}
\makeatletter
\@ifundefined{KOMAClassName}{% if non-KOMA class
  \IfFileExists{parskip.sty}{%
    \usepackage{parskip}
  }{% else
    \setlength{\parindent}{0pt}
    \setlength{\parskip}{6pt plus 2pt minus 1pt}}
}{% if KOMA class
  \KOMAoptions{parskip=half}}
\makeatother
\usepackage{xcolor}
\usepackage[margin=0.8in]{geometry}
\usepackage{longtable,booktabs,array}
\usepackage{calc} % for calculating minipage widths
% Correct order of tables after \paragraph or \subparagraph
\usepackage{etoolbox}
\makeatletter
\patchcmd\longtable{\par}{\if@noskipsec\mbox{}\fi\par}{}{}
\makeatother
% Allow footnotes in longtable head/foot
\IfFileExists{footnotehyper.sty}{\usepackage{footnotehyper}}{\usepackage{footnote}}
\makesavenoteenv{longtable}
\setlength{\emergencystretch}{3em} % prevent overfull lines
\providecommand{\tightlist}{%
  \setlength{\itemsep}{0pt}\setlength{\parskip}{0pt}}
\setcounter{secnumdepth}{-\maxdimen} % remove section numbering
\ifLuaTeX
\usepackage[bidi=basic]{babel}
\else
\usepackage[bidi=default]{babel}
\fi
\babelprovide[main,import]{american}
\ifPDFTeX
\else
\babelfont{rm}[]{Liberation Serif}
\fi
% get rid of language-specific shorthands (see #6817):
\let\LanguageShortHands\languageshorthands
\def\languageshorthands#1{}
\usepackage{adjustbox}
\emergencystretch=3em
\AtBeginEnvironment{longtable}{\small\sloppy}
\usepackage{stmaryrd}
% Printing shims for run-on conditional-bar tokens in the supplied Biology source.
\providecommand{\midL}{\mid L}
\providecommand{\midN}{\mid N}
\sloppy
\ifLuaTeX
  \usepackage{selnolig}  % disable illegal ligatures
\fi
\usepackage{bookmark}
\IfFileExists{xurl.sty}{\usepackage{xurl}}{} % add URL line breaks if available
\urlstyle{same}
\hypersetup{
  pdflang={en-US},
  colorlinks=true,
  linkcolor={blue},
  filecolor={Maroon},
  citecolor={Blue},
  urlcolor={blue},
  pdfcreator={LaTeX via pandoc}}

\author{}
\date{}

\begin{document}

\section{COHERENCE UNIVERSALISM: SOCIAL
DYNAMICS}\label{coherence-universalism-social-dynamics}

\subsection{Coordination, Power, Meaning, and Repair in Plural
Societies}\label{coordination-power-meaning-and-repair-in-plural-societies}

\textbf{Gaura Kiśora Dās Rader}\\
\textbf{Heaven≡Earth Foundation}\\
\textbf{White Paper Series}

\textbf{Version 0.10}\\
\textbf{October 2026}

\textbf{Release status.} This manuscript is part of the complete Version
0.10 proposal. Numbered principles identify the framework's current
formulations; their claim status, external validation, and authorization
for consequential application remain separate questions.

\textbf{Freeze status.} This manuscript is synchronized to the frozen
Version 0.10 release registry, evidence package, and review gates.
Future changes require a new versioned release.

\begin{center}\rule{0.5\linewidth}{0.5pt}\end{center}

\section{Abstract}\label{abstract}

Collective life is organized through persons, relationships, norms,
institutions, infrastructures, technologies, economies, cultures, and
ecosystems. These arrangements make coordinated action possible, but
they also distribute power, risk, recognition, knowledge, opportunity,
and material security. Coherence Universalism describes the social
domain as \textbf{distributed coordination and correction under plural
agency}. That description becomes analytically useful only when the
collective unit, affected groups, authority relations, protected
conditions, observation channels, and temporal horizons are declared.

This paper develops a non-scalar architecture for social analysis. A
social model specifies agents or groups, state variables, environment,
admissible actions, transition processes, protected conditions,
interaction networks, authority relations, distributions of resources
and burdens, public meanings, history, scale, horizon, and observation
model. Social organization is evaluated through a profile rather than
one universal score. Candidate dimensions include material security,
calibrated trust, information integrity, coordination capacity,
meaningful voice, legitimacy, distribution, adaptability, repair
capacity, extraction, coercion, and ecological or intergenerational
burden. These dimensions can conflict. A regime can coordinate
effectively while suppressing voice. A plural institution can protect
dissent while failing to decide. High average welfare can coexist with
severe harm to a minority. Stable compliance can arise from consent,
habit, dependency, or force.

The paper therefore distinguishes coordination from consensus, trust
from trustworthiness, compliance from consent, sociological acceptance
from normative legitimacy, pluralism from fragmentation, and stability
from justice. It treats protected dissent as one possible correction
channel while rejecting the inference that dissent is automatically
accurate. It reconstructs institutions as durable systems of roles,
rules, records, resources, sanctions, memory, and appeal whose
legitimacy depends not only on performance but also on agency,
contestability, burden distribution, independent evidence, and practical
repair.

A deterministic contrast model illustrates the architecture's
non-entailments. In the declared construction, centralized suppression
responds fastest to a clear common emergency, immediate-feedback
contestable polycentricity corrects fastest when a locally exposed group
detects a distributed shift, and majoritarian openness best filters a
one-group false alarm. Three nonnegative scalarizations rank three
different regimes first. A separate delay sweep shows that the adaptive
regime's correction-speed advantage depends on timely and attributable
outcome feedback. These are designed properties of a pedagogical model,
not evidence for one universally preferred political order.

The paper then extends the account of culture into \textbf{Open Myth}: a
shared narrative horizon that can orient identity, responsibility,
sacrifice, and intergenerational action without claiming sovereign
authority over description, identity, or belonging. Myth is treated as
social infrastructure and social power. It can coordinate care or
domination. A legitimate civilizational narrative must therefore remain
truth-constrained, plural in thick expression, participatory in
authorship, contestable, revisable, open to refusal, and unable to
override protected standing or correctability. \emph{The Book of Eight
Gates} is treated as one rich mythic expression of the CU horizon, not
as a mandatory cosmology or scientific authority.

The paper closes with \textbf{CU-COM11 --- Fellowship, Justice, and
Collective Repair}. Social coherence is not conformity or the absence of
conflict. The commitment is to forms of common life in which plural
agency, material capability, protected dissent, truthful public records,
legitimate authority, and practical correction can coexist.

\textbf{Keywords:} social coordination; institutions; power; legitimacy;
dissent; political economy; culture; myth; collective repair; artificial
intelligence

\subsection{Relation to the series}\label{relation-to-the-series}

Psychology shows how meaning, identity, self-deception, trauma, agency,
and repair are organized in persons and relationships. Social Dynamics
begins where these processes are distributed across institutions, norms,
infrastructures, economies, publics, and histories that no individual
controls. Ethics asks how agents and institutions remain truthful and
answerable within this field. Artificial Intelligence extends the
analysis to sociotechnical systems whose models, operators, owners,
users, affected communities, and possible artificial subjects cannot be
treated as one agent.

This paper can be read independently. Its
\href{../supplements/07_Social_Dynamics_Technical_Supplement_v0.10.md}{technical
supplement} preserves the complete formal propositions, deterministic
model, claim ledger, principle registry, independent audit, bounded
repair, literature ledger, and Open Myth source architecture.

\subsection{Note on Version 0.10}\label{note-on-version-0.10}

Version 0.10 retains the original concern with institutions, culture,
civilizational transition, meaning, technology, and collective repair
while replacing a universal social scalar, general coherence gradient,
three-field ontology, universal collapse threshold, and high-coherence
attractor with declared social models, plural profiles, protected
conditions, explicit power analysis, and conditional institutional
hypotheses. The main text presents the current account directly;
Appendix B and the technical archive preserve the revision
record.\footnote{Earlier versions sometimes treated social coherence as
  a scalar and narrated human history through one civilizational
  trajectory. Version 0.10 retains the integrative questions while
  distinguishing coordination, legitimacy, distribution, power,
  correctability, and commitment.}

\textbf{Release-candidate synchronization.} This manuscript uses the
shared Version 0.10 definitions, principle identifiers, Mandala
cross-references, and interpretive and commitment language published in
\emph{Coherence Universalism: Metaphysics and Epistemology} and
\emph{Coherence Universalism: Foundations}. Its audited technical core,
formal results, numerical results, and claim-status boundaries remain
unchanged except where this manuscript explicitly states otherwise.

\subsection{How claims are marked}\label{how-claims-are-marked}

The paper uses \textbf{definition}, \textbf{formal result},
\textbf{imported result}, \textbf{empirical generalization},
\textbf{conditional result}, \textbf{bridge principle},
\textbf{normative principle}, \textbf{hypothesis},
\textbf{interpretation}, \textbf{commitment}, and \textbf{open problem}
in scoped senses. A domain authority identifies the controlling CU
formulation; it does not convert a provisional empirical or normative
claim into established fact.

\section{Contents}\label{contents}

\begin{itemize}
\tightlist
\item
  \begin{enumerate}
  \def\labelenumi{\arabic{enumi}.}
  \tightlist
  \item
    Social Organization Under Plural Agency
  \end{enumerate}
\item
  \begin{enumerate}
  \def\labelenumi{\arabic{enumi}.}
  \setcounter{enumi}{1}
  \tightlist
  \item
    Declaring a Social Model
  \end{enumerate}
\item
  \begin{enumerate}
  \def\labelenumi{\arabic{enumi}.}
  \setcounter{enumi}{2}
  \tightlist
  \item
    Coordination Without Uniformity
  \end{enumerate}
\item
  \begin{enumerate}
  \def\labelenumi{\arabic{enumi}.}
  \setcounter{enumi}{3}
  \tightlist
  \item
    Institutions: Continuity, Legitimacy, and Correctability
  \end{enumerate}
\item
  \begin{enumerate}
  \def\labelenumi{\arabic{enumi}.}
  \setcounter{enumi}{4}
  \tightlist
  \item
    Power, Material Security, and Extraction
  \end{enumerate}
\item
  \begin{enumerate}
  \def\labelenumi{\arabic{enumi}.}
  \setcounter{enumi}{5}
  \tightlist
  \item
    Information Integrity, Culture, and Public Meaning
  \end{enumerate}
\item
  \begin{enumerate}
  \def\labelenumi{\arabic{enumi}.}
  \setcounter{enumi}{6}
  \tightlist
  \item
    Polarization, Conflict, and Democratic Correction
  \end{enumerate}
\item
  \begin{enumerate}
  \def\labelenumi{\arabic{enumi}.}
  \setcounter{enumi}{7}
  \tightlist
  \item
    Meso-Scale Social Infrastructure
  \end{enumerate}
\item
  \begin{enumerate}
  \def\labelenumi{\arabic{enumi}.}
  \setcounter{enumi}{8}
  \tightlist
  \item
    Crisis, Cascades, Resilience, and Thresholds
  \end{enumerate}
\item
  \begin{enumerate}
  \def\labelenumi{\arabic{enumi}.}
  \setcounter{enumi}{9}
  \tightlist
  \item
    Social Repair
  \end{enumerate}
\item
  \begin{enumerate}
  \def\labelenumi{\arabic{enumi}.}
  \setcounter{enumi}{10}
  \tightlist
  \item
    Technology, Platforms, and AI-Mediated Social Systems
  \end{enumerate}
\item
  \begin{enumerate}
  \def\labelenumi{\arabic{enumi}.}
  \setcounter{enumi}{11}
  \tightlist
  \item
    Historical and Civilizational Interpretation
  \end{enumerate}
\item
  \begin{enumerate}
  \def\labelenumi{\arabic{enumi}.}
  \setcounter{enumi}{12}
  \tightlist
  \item
    Myth, Narrative, and Civilizational Alignment
  \end{enumerate}
\item
  \begin{enumerate}
  \def\labelenumi{\arabic{enumi}.}
  \setcounter{enumi}{13}
  \tightlist
  \item
    Research Program and Principle Architecture
  \end{enumerate}
\item
  \begin{enumerate}
  \def\labelenumi{\arabic{enumi}.}
  \setcounter{enumi}{14}
  \tightlist
  \item
    What Commitments This Account Licenses
  \end{enumerate}
\item
  \begin{enumerate}
  \def\labelenumi{\arabic{enumi}.}
  \setcounter{enumi}{15}
  \tightlist
  \item
    Conclusion: A Common World Without Enforced Sameness
  \end{enumerate}
\item
  Appendix A. Social Principles and Claim-Governance Summary
\item
  Appendix B. Major Revisions from v0.09
\item
  Appendix C. Technical Supplement and Reproducibility
\item
  References
\end{itemize}

\section{Principal Contributions}\label{principal-contributions}

\begin{enumerate}
\def\labelenumi{\arabic{enumi}.}
\tightlist
\item
  Defines social organization as distributed coordination and correction
  under plural agency.
\item
  Replaces one social-coherence score with profiles, protected floors,
  distributional reporting, and explicit scalarization choices.
\item
  Makes power, authority, agenda setting, ownership, sanctions, and exit
  costs constitutive variables rather than afterthoughts.
\item
  Distinguishes coordination, compliance, consent, trust,
  trustworthiness, legitimacy, and justice.
\item
  Defines practical correctability through accessible and sufficiently
  independent challenge, reversal, remedy, recurrence tracking, and
  learning.
\item
  Treats dissent as a protected but fallible correction channel.
\item
  Integrates material security, labor, extraction, attention, ecological
  burden, and future generations into social analysis.
\item
  Develops Open Myth and Addressable Myth as truth-constrained forms of
  narrative orientation under plural authorship and exit.
\item
  Locates artificial intelligence inside social and civilizational
  governance rather than treating it as a merely technical object.
\item
  Articulates fellowship, justice, and collective repair as the Social
  Dynamics commitment principle.
\end{enumerate}

\begin{center}\rule{0.5\linewidth}{0.5pt}\end{center}

\section{Domain Analytic Frame and CAP
Record}\label{domain-analytic-frame-and-cap-record}

\textbf{Target:} Coherence Universalism Version 0.10 release\\
\textbf{Status:} Release declaration; scope-limited and revisable under
the declared review gates\\
\textbf{Use rule:} Unresolved fields are left explicit rather than
completed by analogy.

\subsection{Paper and model identity}\label{paper-and-model-identity}

\begin{longtable}[]{@{}
  >{\raggedright\arraybackslash}p{(\columnwidth - 4\tabcolsep) * \real{0.3333}}
  >{\raggedright\arraybackslash}p{(\columnwidth - 4\tabcolsep) * \real{0.3333}}
  >{\raggedright\arraybackslash}p{(\columnwidth - 4\tabcolsep) * \real{0.3333}}@{}}
\toprule\noalign{}
\begin{minipage}[b]{\linewidth}\raggedright
Field
\end{minipage} & \begin{minipage}[b]{\linewidth}\raggedright
Declaration
\end{minipage} & \begin{minipage}[b]{\linewidth}\raggedright
Status
\end{minipage} \\
\midrule\noalign{}
\endhead
\bottomrule\noalign{}
\endlastfoot
Paper & Coherence Universalism: Social Dynamics & Specified \\
Domain owner & Social Dynamics & Specified by the Version 0.10 release
corpus \\
Primary model or analysis & Social model
\(\mathcal S=(A,X,E,U,F,K,G,P,R,N,\mathcal H,\pi,T,\mathcal O)\) &
Paper-relative \\
Primary claim status & Definitions, empirical generalizations,
diagnostic principles, conditional models, normative bridges,
interpretations, and commitments & Mixed statuses remain explicit \\
Review gate & Social science, political economy, institutional design,
history, law, ecology, measurement, and affected-community review &
Proportionate to claim and consequential reach \\
Methodological authority & CU-FP2, CU-M2, CU-M11, CU-M12, CU-M13,
CU-M14, and CU-M15 & Controlling Version 0.10 methodological
framework \\
\end{longtable}

\subsection{Shared Analytic Frame}\label{shared-analytic-frame}

\begin{longtable}[]{@{}
  >{\raggedright\arraybackslash}p{(\columnwidth - 4\tabcolsep) * \real{0.3333}}
  >{\raggedright\arraybackslash}p{(\columnwidth - 4\tabcolsep) * \real{0.3333}}
  >{\raggedright\arraybackslash}p{(\columnwidth - 4\tabcolsep) * \real{0.3333}}@{}}
\toprule\noalign{}
\begin{minipage}[b]{\linewidth}\raggedright
\#
\end{minipage} & \begin{minipage}[b]{\linewidth}\raggedright
Slot and domain declaration
\end{minipage} & \begin{minipage}[b]{\linewidth}\raggedright
Status and evidence
\end{minipage} \\
\midrule\noalign{}
\endhead
\bottomrule\noalign{}
\endlastfoot
1 & \textbf{Unit and boundary.} A declared collective, institution,
network, jurisdiction, public, economy, or sociotechnical field; society
is never an unmarked totality. & Boundary and jurisdiction required \\
2 & \textbf{State or trajectory space.} \(X\) plus agents, institutions,
resources, norms, histories, and distributions represent social states
and trajectories. & Mixed quantitative and qualitative \\
3 & \textbf{Environment.} \(E\) includes ecological, technological,
geopolitical, material, and cross-jurisdictional conditions. &
Explicit \\
4 & \textbf{Constraints and viable region.} \(K\) contains protected
standing, material floors, agency, dissent, legitimacy, ecological, and
intergenerational conditions. & Normative and empirical components
distinguished \\
5 & \textbf{Actions, options, and transitions.} \(U\) contains policy,
institutional response, collective action, refusal, exit, enforcement,
bargaining, and repair; option generation is power-laden. & Declared \\
6 & \textbf{Dynamics or update rule.} \(F\) may include strategic
action, learning, bargaining, enforcement, drift, shock response, and
feedback; no universal social gradient is claimed. & Model-specific \\
7 & \textbf{Observation and access.} \(\mathcal O\) includes surveys,
records, indicators, testimony, public data, and feedback institutions
with omission and delay. & Measurement institution must be declared \\
8 & \textbf{Perturbation and intervention class.} Policy change, shock,
conflict, information intervention, institutional reform, material
support, and governance redesign are distinct classes. &
Application-specific \\
9 & \textbf{Equilibria, basins, or regimes.} The paper compares regimes
in constructed models but does not demonstrate one high-coherence
attractor or universal basin structure. & Regime language only where
supported \\
10 & \textbf{Local--global status.} Distribution matrices, subgroup
effects, cross-scale externalization, and delayed burdens are primary.
Exact decompositions are possible only for separately licensed scalar
distributions. & Coordinate-wise / qualitative by default \\
11 & \textbf{Scalar status.} The social object is a twelve-coordinate
profile. Decision scalarizations may rank admissible options for a
declared purpose but not societies in general. & No universal social
scalar \\
12 & \textbf{Statistical-image and allocation status.} Population
distributions and genuine shares of income, resources, exposure, voice,
or attention may qualify. Trust, legitimacy, and agency profiles do not.
& Strong allocation candidates plus non-simplex profiles \\
13 & \textbf{Translation status.} Institutional, capability, trust,
information, network, and inequality formalisms transfer only with their
assumptions and target validation. & Levels 1--4 \\
14 & \textbf{Open gaps.} Measurement of correctability, power,
legitimacy, extraction, feedback quality, option adequacy, and
cross-scale repair remains incomplete. & Open \\
\end{longtable}

\subsection{Compact Coherence Analysis Protocol
record}\label{compact-coherence-analysis-protocol-record}

\begin{longtable}[]{@{}
  >{\raggedright\arraybackslash}p{(\columnwidth - 4\tabcolsep) * \real{0.3333}}
  >{\raggedright\arraybackslash}p{(\columnwidth - 4\tabcolsep) * \real{0.3333}}
  >{\raggedright\arraybackslash}p{(\columnwidth - 4\tabcolsep) * \real{0.3333}}@{}}
\toprule\noalign{}
\begin{minipage}[b]{\linewidth}\raggedright
CAP module
\end{minipage} & \begin{minipage}[b]{\linewidth}\raggedright
Declaration
\end{minipage} & \begin{minipage}[b]{\linewidth}\raggedright
Boundary
\end{minipage} \\
\midrule\noalign{}
\endhead
\bottomrule\noalign{}
\endlastfoot
System declaration & Name agents and groups, boundary, jurisdiction,
environment, authority, distributions, norms, history, and horizon. &
Section 2.1 \\
Observation declaration & Declare records, categories, sampled
populations, missing groups, feedback delay, attribution, and control of
data. & Aggregation can hide harm \\
Evaluation declaration & Use profiles and distribution matrices, declare
floors and option sets, preserve non-dominance, and state any scalar
purpose. & Distribution before aggregation \\
Power and authorship declaration & Name agenda setters, rule makers,
owners, classifiers, sanctioning authorities, beneficiaries, burden
bearers, and routes of challenge or exit. & Power is a required
variable \\
Revision declaration & Changes to affected groups, categories,
feasibility, jurisdiction, floor, horizon, or authority are model
revision. & Must be auditable \\
\end{longtable}

\subsection{License summary}\label{license-summary}

\begin{longtable}[]{@{}
  >{\raggedright\arraybackslash}p{(\columnwidth - 4\tabcolsep) * \real{0.3333}}
  >{\raggedright\arraybackslash}p{(\columnwidth - 4\tabcolsep) * \real{0.3333}}
  >{\raggedright\arraybackslash}p{(\columnwidth - 4\tabcolsep) * \real{0.3333}}@{}}
\toprule\noalign{}
\begin{minipage}[b]{\linewidth}\raggedright
License
\end{minipage} & \begin{minipage}[b]{\linewidth}\raggedright
Declaration
\end{minipage} & \begin{minipage}[b]{\linewidth}\raggedright
Current status
\end{minipage} \\
\midrule\noalign{}
\endhead
\bottomrule\noalign{}
\endlastfoot
Scalar license & Coordinate measures and declared decision
scalarizations; no total social score. & Conditional \\
Local--global license & Exact only for validated decomposable
indicators; otherwise coordinate-wise or qualitative. & Mixed \\
Statistical Image License & Licensed for declared population, outcome,
exposure, and opinion distributions. & Object-specific \\
Allocation License & Strong where a real common-unit total is divided
among groups, institutions, or uses. & Licensed when the total and
partition are real \\
Refinement and Coarse-Graining License & Permitted where group
aggregation, decomposition, and category changes are explicit and do not
erase affected parties. & Conditional \\
Fisher-Rao License & Possible for licensed distributions and
allocations; it does not measure legitimacy, trustworthiness, or
justice. & Object-specific only \\
Translation license & Cross-field methods require both structural fit
and target-domain validation. & Levels 2--4 as earned \\
\end{longtable}

\subsection{Open gaps and review
disposition}\label{open-gaps-and-review-disposition}

\begin{longtable}[]{@{}
  >{\raggedright\arraybackslash}p{(\columnwidth - 4\tabcolsep) * \real{0.3333}}
  >{\raggedright\arraybackslash}p{(\columnwidth - 4\tabcolsep) * \real{0.3333}}
  >{\raggedright\arraybackslash}p{(\columnwidth - 4\tabcolsep) * \real{0.3333}}@{}}
\toprule\noalign{}
\begin{minipage}[b]{\linewidth}\raggedright
Gap
\end{minipage} & \begin{minipage}[b]{\linewidth}\raggedright
Consequence
\end{minipage} & \begin{minipage}[b]{\linewidth}\raggedright
Required work
\end{minipage} \\
\midrule\noalign{}
\endhead
\bottomrule\noalign{}
\endlastfoot
Distributional measurement & Means can conceal floor violations and
subgroup deterioration. & Disaggregated longitudinal data \\
Option-set adequacy & A valid ranking can operate on a captured menu. &
Formal and participatory adequacy tests \\
Institutional correctability & Complaint channels may be symbolic. &
Independent reversal, remedy, and recurrence evidence \\
\end{longtable}

\begin{quote}
\textbf{Propagation rule.} This card makes the paper's existing objects
and limits visible. It does not strengthen any empirical, mathematical,
normative, metaphysical, or spiritual claim beyond the wording supported
in the manuscript.
\end{quote}

\begin{center}\rule{0.5\linewidth}{0.5pt}\end{center}

\section{1. Social organization under plural
agency}\label{social-organization-under-plural-agency}

\subsection{1.1 Why collective life is a distinct
domain}\label{why-collective-life-is-a-distinct-domain}

A society is not one mind enlarged. It is not an organism in any
straightforward biological sense, and it is not a single optimizer whose
purposes can be read from aggregate behavior. It consists of partially
autonomous and unequally situated agents whose actions are linked by
dependence, communication, rules, infrastructure, history, and power.

Yet social life is more than an accidental aggregate. Language lets
expectations travel between persons. Norms make some actions
intelligible and others costly. Institutions preserve records and roles
beyond individual memory. Markets coordinate dispersed exchanges.
Families, congregations, neighborhoods, schools, workplaces,
professions, movements, and states create durable relations among people
who cannot continually renegotiate every condition of their common life.
Culture carries practices and meanings across generations. Material
systems make coordination possible over distance and time.

\textbf{{[}Definition{]} Social organization} is the historically
extended pattern of relations, rules, expectations, resources,
infrastructures, and meaning structures through which partially
autonomous agents coordinate, conflict, cooperate, exercise power, and
respond to shared or imposed conditions.

This definition is broader than social harmony. Conflict can be
organized. Domination can be organized. Resistance can be organized.
Social analysis must explain all of them.

\subsection{1.2 The governing question}\label{the-governing-question}

The central social question is not simply why societies hold together.
Enslaving institutions, colonial administrations, caste orders, criminal
networks, authoritarian states, and exploitative firms can be durable
and internally coordinated. Persistence does not convert organization
into legitimacy.

Three questions must therefore remain separate:

\begin{enumerate}
\def\labelenumi{\arabic{enumi}.}
\item
  \textbf{How is coordination produced?}\\
  Which rules, incentives, conventions, infrastructures, identities, and
  authority relations generate mutually adjusted action?
\item
  \textbf{How are benefits, burdens, voice, and risk distributed?}\\
  Who gains capacity, who loses it, and whose experience disappears
  inside an aggregate?
\item
  \textbf{Under what conditions is coordination legitimate and
  correctable?}\\
  Can those subject to a system understand, contest, revise, leave, or
  repair it without catastrophic retaliation or loss of basic standing?
\end{enumerate}

A complete account of social organization must address all three.

\subsection{1.3 Social coherence as a bridge
concept}\label{social-coherence-as-a-bridge-concept}

\textbf{{[}Bridge principle{]} Social coherence} refers to integrated
social organization under declared constraints: the capacity of
partially autonomous agents and institutions to coordinate, contest,
correct, and materially support collective life under disagreement and
change.

The term is useful only if the relevant unit, protected conditions,
affected groups, power relations, and horizon are stated. A family,
city, scientific community, platform, labor market, state, and
transnational supply chain are not interchangeable units. A system can
be coherent for one constituency and destructive for another.

Social coherence is therefore not:

\begin{itemize}
\tightlist
\item
  harmony;
\item
  consensus;
\item
  conformity;
\item
  institutional persistence;
\item
  majority satisfaction;
\item
  aggregate wealth;
\item
  or moral goodness.
\end{itemize}

It is a relational description that must be joined to legitimacy,
distribution, and repair.

\subsection{1.4 Why the social domain cannot be reduced to
psychology}\label{why-the-social-domain-cannot-be-reduced-to-psychology}

Individual beliefs and emotions matter, but social outcomes are also
produced by structures no individual controls. A worker may distrust a
firm because wages are insecure, not because of an attachment style. A
population may show anxiety because housing, debt, care, or ecological
conditions are unstable. A public may polarize because institutions
reward identity sorting and strategic amplification. A person may comply
with a rule because exit is impossible.

Psychological processes and social structures interact. Neither
substitutes for the other. Social explanation begins when roles,
incentives, records, law, infrastructure, property, public meaning, and
unequal power make a causal difference beyond individual dispositions.

\subsection{1.5 Conceptual prohibitions}\label{conceptual-prohibitions}

Every later claim is constrained by the following rules:

\begin{enumerate}
\def\labelenumi{\arabic{enumi}.}
\tightlist
\item
  \textbf{No scalar by default.} Social profiles may be scalarized only
  for a declared purpose with disclosed weights, normalizations, floors,
  affected groups, and horizon.
\item
  \textbf{No legitimacy from order.} Coordination, compliance,
  popularity, and persistence do not establish consent or justice.
\item
  \textbf{No average without distribution.} Aggregate gains may conceal
  concentrated or noncompensable harm.
\item
  \textbf{No trust without trustworthiness.} Reliance can be misplaced,
  coerced, or strategically produced.
\item
  \textbf{No dissent romanticism.} Dissent must be protected without
  being presumed correct.
\item
  \textbf{No institution as its own sole auditor.} Nominal complaint
  systems and self-reported success do not establish practical
  correctability.
\item
  \textbf{No politics by clinical diagnosis.} Individual psychological
  constructs may not substitute for institutional, economic, historical,
  or political explanation.
\item
  \textbf{No mythic sovereignty.} Shared narrative may orient common
  life but may not override truth, standing, dissent, refusal, or
  correction.
\end{enumerate}

\subsection{1.6 Roadmap}\label{roadmap}

Sections 2--5 develop the model, coordination, institutional, and
political-economic architecture. Sections 6--8 examine information,
culture, conflict, and meso-scale infrastructure. Sections 9--11 address
crisis, repair, and technology. Section 12 gives the historical
interpretation. Section 13 develops Open Myth and civilizational
alignment. Section 14 presents the research program and principle
architecture. Section 15 states the licensed commitments. Section 16
gathers the account into a bounded conclusion.

\begin{center}\rule{0.5\linewidth}{0.5pt}\end{center}

\section{2. Declaring a social model}\label{declaring-a-social-model}

\subsection{2.1 The model tuple}\label{the-model-tuple}

A candidate social model is represented as

\[\adjustbox{max width=\linewidth}{$\displaystyle 
\mathcal S
=
(A,X,E,U,F,K,G,P,R,N,\mathcal H,\pi,T,\mathcal O).
$}\]

The components are:

\begin{itemize}
\tightlist
\item
  \(A\): agents, groups, organizations, and other affected parties;
\item
  \(X\): relevant social and material state variables;
\item
  \(E\): environmental, ecological, technological, and geopolitical
  conditions;
\item
  \(U\): admissible actions, rules, policies, and institutional
  responses;
\item
  \(F\): transition processes, including strategic action, learning,
  bargaining, enforcement, drift, and shock response;
\item
  \(K\): protected constraints or floors;
\item
  \(G\): interaction, communication, dependency, and exchange relations;
\item
  \(P\): authority, agenda-setting, sanctioning, ownership, and
  epistemic power;
\item
  \(R\): resource, benefit, risk, and burden distributions;
\item
  \(N\): norms, categories, narratives, identities, and public meaning
  structures;
\item
  \(\mathcal H\): relevant history and path dependence;
\item
  \(\pi\): boundary, scale, jurisdiction, and access map;
\item
  \(T\): temporal horizon or set of horizons; and
\item
  \(\mathcal O\): observation and measurement channels.
\end{itemize}

Not every application requires a fully quantified object. The tuple is a
declaration discipline. It prevents an analyst from speaking of
``society'' while measuring one population, one platform, one nation, or
one survey response.

The standardized Domain Analytic Frame above is the compact CAP record
for this paper. It makes the paper's scalar, local--global, allocation,
and translation status visible without replacing the stronger
social-governance rules below.

\subsection{2.2 A plural social profile}\label{a-plural-social-profile}

This paper replaces the retired scalar with the profile

\[\adjustbox{max width=\linewidth}{$\displaystyle 
\mathbf P_{\mathrm{soc}}
=
\langle
p_{\mathrm{mat}},p_{\mathrm{trust}},p_{\mathrm{info}},p_{\mathrm{coord}},p_{\mathrm{voice}},p_{\mathrm{leg}},p_{\mathrm{dist}},p_{\mathrm{adapt}},p_{\mathrm{repair}},p_{\mathrm{extract}},p_{\mathrm{coerce}},p_{\mathrm{eco}}
\rangle.
$}\]

The coordinates are candidate dimensions, not universally fixed
measures:

\begin{itemize}
\tightlist
\item
  \(p_{\mathrm{mat}}\): material security and access to basic
  capabilities;
\item
  \(p_{\mathrm{trust}}\): calibrated trust and demonstrated
  trustworthiness;
\item
  \(p_{\mathrm{info}}\): information integrity, provenance, and
  correction;
\item
  \(p_{\mathrm{coord}}\): coordination capacity;
\item
  \(p_{\mathrm{voice}}\): voice, contestability, refusal, and meaningful
  participation;
\item
  \(p_{\mathrm{leg}}\): legitimacy conditions;
\item
  \(p_{\mathrm{dist}}\): distribution of benefit, burden, standing, and
  risk;
\item
  \(p_{\mathrm{adapt}}\): adaptability and learning;
\item
  \(p_{\mathrm{repair}}\): repair capacity;
\item
  \(p_{\mathrm{extract}}\): extraction and externalization;
\item
  \(p_{\mathrm{coerce}}\): coercion and domination; and
\item
  \(p_{\mathrm{eco}}\): ecological and intergenerational burden.
\end{itemize}

The first nine are ordinarily oriented so that more is preferable only
after their content is specified. The final three ordinarily record
costs or constraints. Even this orientation can require qualification.
Trust in a deceptive source is not an improvement. Rapid adaptability
can become policy volatility. Strong collective memory can preserve
domination as well as wisdom.

A profile may support:

\begin{itemize}
\tightlist
\item
  Pareto comparison;
\item
  threshold or floor analysis;
\item
  lexicographic priority;
\item
  scenario comparison;
\item
  partial ordering; or
\item
  a declared scalarization for a particular decision.
\end{itemize}

It does not automatically support a total ranking of societies.

\subsection{2.3 Distribution before
aggregation}\label{distribution-before-aggregation}

Let \(g\in A\) index affected groups and \(j\) index dimensions. The
primary evaluative object is often a matrix

\[\adjustbox{max width=\linewidth}{$\displaystyle 
Y=(y_{gj}),
$}\]

not a population average.

An average can rise while a protected group falls below a floor.
Formally, there exist profiles \(a,b\) and a floor \(k\) such that

\[\adjustbox{max width=\linewidth}{$\displaystyle 
\frac{1}{|A|}\sum_g a_g
>
\frac{1}{|A|}\sum_g b_g,
$}\]

while

\[\adjustbox{max width=\linewidth}{$\displaystyle 
\min_g a_g<k
\quad\text{and}\quad
\min_g b_g\ge k.
$}\]

For example,

\[\adjustbox{max width=\linewidth}{$\displaystyle 
a=(0.90,0.90,0.20),
\qquad
b=(0.65,0.65,0.65),
\qquad
k=0.40.
$}\]

The first arrangement has the higher mean and the unacceptable floor
violation. This is not an argument that averages are useless. It is an
argument that distribution and noncompensable conditions cannot be
reconstructed from averages after the fact.

\subsection{2.4 Power as a required
variable}\label{power-as-a-required-variable}

\textbf{{[}Definition{]} Social power} is a relational capacity to shape
another's available actions, costs, information, categories,
opportunities, standing, or exposure to consequences.

Power includes more than direct command. It can operate through:

\begin{itemize}
\tightlist
\item
  ownership and control of material resources;
\item
  rule making and enforcement;
\item
  agenda setting;
\item
  control of communication channels;
\item
  classification and measurement;
\item
  institutional role;
\item
  credentialing and expertise;
\item
  dependency and exit costs;
\item
  network position;
\item
  threat, violence, or surveillance;
\item
  and the ability to define which experiences are publicly intelligible.
\end{itemize}

A social profile that records trust and coordination but omits power can
mistake successful domination for voluntary integration.

\subsection{2.5 Protected conditions}\label{protected-conditions}

Protected conditions are not ordinary weighted preferences. Depending on
the domain, they may include:

\begin{itemize}
\tightlist
\item
  bodily safety;
\item
  basic subsistence;
\item
  equal legal standing;
\item
  freedom from torture, enslavement, and arbitrary detention;
\item
  meaningful consent;
\item
  access to appeal;
\item
  non-retaliatory channels for reporting harm;
\item
  domains of privacy and association;
\item
  ecological conditions necessary for continued life;
\item
  and minimum agency for those subject to institutional power.
\end{itemize}

The content of these floors is normative and political. Mathematics does
not select it. The model must disclose who chose the floors, how
affected groups participated, and what happens when protections
conflict.

\subsection{2.6 Trust, trustworthiness, and
reliance}\label{trust-trustworthiness-and-reliance}

Trust is often treated as an uncomplicated social good. It is not.

\textbf{Trust} is a disposition to accept vulnerability on the
expectation that another person or institution will act with sufficient
competence, reliability, and regard. \textbf{Trustworthiness} concerns
whether that expectation is warranted. \textbf{Reliance} can occur
without trust, as when a person uses a system because no alternative
exists.

Accordingly:

\begin{itemize}
\tightlist
\item
  low trust can reflect paranoia, historical memory, or accurate
  appraisal;
\item
  high trust can reflect deserved confidence, dependency, propaganda,
  habit, or lack of alternatives;
\item
  trust can be domain-specific;
\item
  trust can be asymmetrical; and
\item
  institutions may seek the appearance of trust while resisting the
  conditions that would make them trustworthy.
\end{itemize}

The OECD's cross-national trust work is useful partly because it
disaggregates experiences and expectations concerning reliability,
responsiveness, openness, integrity, and fairness rather than treating
trust as one free-floating attitude (OECD, 2024).

\subsection{2.7 Compliance, consent, and
legitimacy}\label{compliance-consent-and-legitimacy}

Observed compliance is behavior. Consent is a relation between an agent
and a choice under adequate information and voluntariness. Sociological
legitimacy is the belief that an authority is appropriate or entitled to
rule. Normative legitimacy concerns whether the authority can be
justified to those subject to it.

These can diverge.

\[\adjustbox{max width=\linewidth}{$\displaystyle 
\text{compliance}
\not\Rightarrow
\text{consent}
\not\Rightarrow
\text{normative legitimacy}.
$}\]

A population may comply because rules are accepted, because sanctions
are severe, because exit is impossible, because alternatives are
invisible, or because coordination itself is valuable despite
disagreement. Tyler's legitimacy research emphasizes that procedural
treatment and perceived appropriateness can motivate voluntary
cooperation, while also illustrating why legitimacy must not be reduced
to force or reward (Tyler, 2006, 2023). The further point is that
perceived legitimacy is itself not the final normative verdict.

\subsection{2.8 Measurement as
intervention}\label{measurement-as-intervention}

Social indicators do not merely observe. Once tied to funding, status,
punishment, or promotion, they alter behavior. Institutions may improve
a proxy rather than the condition it was meant to represent. Records can
become targets. Categories can determine eligibility. What is legible to
administrators can displace what matters to affected people.

Every proposed social metric should therefore state:

\begin{itemize}
\tightlist
\item
  its construct;
\item
  data source;
\item
  missing groups;
\item
  measurement error;
\item
  incentives it creates;
\item
  likely gaming strategies;
\item
  distributional effects;
\item
  and conditions under which the metric should be retired.
\end{itemize}

This is the core of CU-ICF-1, Metric Absolutism, and CU-ICF-5, Silent
Optimization.

\subsection{2.9 Nineteen governance
rules}\label{nineteen-governance-rules}

Every application of this framework should satisfy the following rules.

\begin{enumerate}
\def\labelenumi{\arabic{enumi}.}
\tightlist
\item
  Name the social unit, jurisdiction, and boundary.
\item
  Identify affected groups, including those outside the formal decision
  body.
\item
  Map authority, ownership, agenda-setting power, sanctions, and exit
  costs.
\item
  State the temporal horizon and any discounting of future burdens.
\item
  Declare protected floors separately from weighted benefits.
\item
  Report distributions, not only averages.
\item
  Distinguish trust from trustworthiness and compelled reliance.
\item
  Distinguish compliance, consent, sociological legitimacy, and
  normative legitimacy.
\item
  Preserve independent evidence channels and disclose provenance.
\item
  Protect dissent without treating dissent as automatically true.
\item
  Distinguish coordination capacity from justice or moral value.
\item
  Do not restore a universal social coherence scalar.
\item
  Declare proxies, missing data, and measurement-induced incentives.
\item
  Test for metric capture, gaming, and classification harm.
\item
  Distinguish association, mechanism, intervention effect, and
  interpretation.
\item
  Do not diagnose populations with individual clinical concepts.
\item
  Give affected communities substantive authority in classifying harms
  and evaluating repair.
\item
  Keep spiritual and civilizational interpretation from substituting for
  empirical or normative argument.
\item
  Do not treat nominal appeals, complaint channels, or self-reported
  review statistics as evidence of correctability without
  practical-access, independence, remedy, recurrence, and affected-party
  evidence.
\end{enumerate}

\section{3. Coordination without
uniformity}\label{coordination-without-uniformity}

\subsection{3.1 Mutual adjustment}\label{mutual-adjustment}

Social coordination occurs when agents make their actions mutually
intelligible and sufficiently compatible. Mechanisms include language,
conventions, prices, schedules, legal rules, professional standards,
rituals, role expectations, reputations, sanctions, shared
infrastructure, and deliberate planning.

No single mechanism is universally best. Prices can coordinate dispersed
exchange while excluding unpriced goods. Hierarchies can act rapidly
while suppressing local knowledge. Deliberation can integrate
perspectives while imposing time and participation burdens. Norms can
reduce transaction costs while preserving unjust expectations. Informal
trust can support cooperation while excluding outsiders.

\textbf{{[}Bridge principle{]} Coordination through mutual adjustment}
is therefore the capacity of agents to revise expectations and action in
response to one another through a declared set of signals, rules, and
feedback channels.

\subsection{3.2 Why consensus is neither necessary nor
sufficient}\label{why-consensus-is-neither-necessary-nor-sufficient}

A coherent disagreement preserves enough shared structure for conflict
to remain corrigible. Participants can identify the issue, distinguish
evidence from loyalty, recognize one another's standing, and use
procedures that do not make defeat equivalent to social annihilation.

Consensus can fail these tests. A unanimous committee may share one
blind spot. A close-knit organization may punish accurate warning. A
public may agree because excluded people are absent from the room. A
doctrine can become internally stable by preventing counterevidence from
entering.

Conversely, visible disagreement may indicate a functioning correction
system. Scientific controversy, adversarial legal procedure,
investigative journalism, labor organizing, protected opposition, and
community protest can increase short-run friction while exposing errors
or burdens hidden by aggregate success.

The task is not to maximize disagreement. It is to organize conflict so
that relevant information and interests can alter decisions without
requiring violence, exile, or total institutional breakdown.

\subsection{3.3 Distributed knowledge}\label{distributed-knowledge}

Knowledge in social systems is dispersed. People occupy different
locations, roles, professions, bodies, and histories. A central
authority may possess broad coordination capacity while lacking local
information. A local group may know its own conditions while lacking
visibility into system-wide effects.

This creates a dual problem:

\[\adjustbox{max width=\linewidth}{$\displaystyle 
\text{local knowledge without common coordination}
\quad\text{and}\quad
\text{central coordination without local knowledge}.
$}\]

Polycentric and federal arrangements are one response. They distribute
decision centers while preserving mechanisms for coordination, learning,
and escalation across scales (Ostrom, 1990, 2010). Such arrangements are
not automatically inclusive or efficient. Local units can reproduce
domination; multiple veto points can block action; coordination costs
can become severe. Their value depends on the problem's scale,
externalities, knowledge distribution, and power structure.

\subsection{3.4 Collective intelligence and its
limits}\label{collective-intelligence-and-its-limits}

Groups can outperform individuals when they combine partially
independent information, relevant diversity, and effective aggregation.
They can also amplify bias, conformity, prestige, or correlated error.
Social influence may improve coordination while reducing the diversity
of estimates that made collective judgment useful.

The concept of collective intelligence therefore requires at least four
questions:

\begin{enumerate}
\def\labelenumi{\arabic{enumi}.}
\tightlist
\item
  Which task is being solved?
\item
  What information is distributed across participants?
\item
  How independent are errors?
\item
  By what process are contributions selected and combined?
\end{enumerate}

A crowd is not wise because it is large. A group is not intelligent
because its output is coherent. The architecture of information and
power matters.

\subsection{3.5 Protected dissent}\label{protected-dissent}

\textbf{{[}Normative bridge principle{]} Protected dissent} is the
institutionally secured capacity to present inconvenient claims,
evidence, objections, or refusals without losing basic standing merely
because they challenge a dominant policy or identity.

Protected dissent serves at least four functions:

\begin{itemize}
\tightlist
\item
  epistemic correction;
\item
  representation of affected interests;
\item
  constraint on power;
\item
  and generation of alternatives.
\end{itemize}

It also has costs. False claims consume attention. Strategic obstruction
can exploit procedural openness. Harassment can masquerade as dissent. A
complete design therefore combines voice with evidence standards,
relevance rules, anti-retaliation protections, transparent adjudication,
and appeal.

The principle is not ``believe the dissenter.'' It is ``do not structure
the institution so that only approved conclusions can reach the place
where correction occurs.''

\subsection{3.6 Majorities, minorities, and
standing}\label{majorities-minorities-and-standing}

Majority decision can provide an intelligible rule for collective
choice, but majority support does not erase minority rights or local
knowledge. The affected-group matrix matters because a policy can be
moderately beneficial to many and catastrophic to a few.

A legitimate majority procedure therefore requires more than counting.
It requires protection of basic standing, disclosure of concentrated
burdens, avenues for review, and rules for cases in which those who
decide do not bear the principal costs.

\subsection{3.7 Coordination under
emergency}\label{coordination-under-emergency}

Emergencies can change the appropriate balance among speed, redundancy,
participation, and discretion. Centralized command may coordinate action
rapidly when the threat is clear, the objective narrow, and delay
dangerous. The same architecture can become maladaptive when uncertainty
is high, local information is decisive, or emergency powers persist
after the emergency.

Accordingly, emergency authority should specify:

\begin{itemize}
\tightlist
\item
  trigger conditions;
\item
  scope;
\item
  duration;
\item
  information channels;
\item
  review;
\item
  termination;
\item
  and repair for wrongful burdens.
\end{itemize}

Temporary coordination need not become permanent sovereignty.

\subsection{3.8 Reader-facing institutional-correction
construction}\label{reader-facing-institutional-correction-construction}

The released technical supplement contains a deterministic three-group,
sixty-step contrast among centralized suppression, majoritarian
openness, and contestable polycentricity. The construction is not an
empirical model of any polity. It is designed to show that speed, local
knowledge, false-alarm filtering, protected voice, burden distribution,
and correction can rank institutional architectures differently.

\begin{longtable}[]{@{}
  >{\raggedright\arraybackslash}p{(\columnwidth - 6\tabcolsep) * \real{0.2000}}
  >{\raggedleft\arraybackslash}p{(\columnwidth - 6\tabcolsep) * \real{0.2667}}
  >{\raggedleft\arraybackslash}p{(\columnwidth - 6\tabcolsep) * \real{0.2667}}
  >{\raggedleft\arraybackslash}p{(\columnwidth - 6\tabcolsep) * \real{0.2667}}@{}}
\toprule\noalign{}
\begin{minipage}[b]{\linewidth}\raggedright
Context and metric
\end{minipage} & \begin{minipage}[b]{\linewidth}\raggedleft
Centralized
\end{minipage} & \begin{minipage}[b]{\linewidth}\raggedleft
Majoritarian
\end{minipage} & \begin{minipage}[b]{\linewidth}\raggedleft
Polycentric
\end{minipage} \\
\midrule\noalign{}
\endhead
\bottomrule\noalign{}
\endlastfoot
Distributed shift: six-step mismatch & 2.259473 & 1.926166 & 1.112662 \\
Distributed shift: correction time & Not reached & 13 steps & 6 steps \\
Clear emergency: correction time & 3 steps & 4 steps & 5 steps \\
One-group false alarm: cumulative deviation & 0.053254 & 0.000000 &
0.167381 \\
\end{longtable}

In the distributed-shift case, the locally exposed group sees the change
first. The polycentric architecture initially corrects fastest because
it includes that signal and receives per-group outcome feedback; the
centralized case gives the exposed group the least voice and highest
burden and violates the declared material and voice floors. In the clear
common emergency, all signals agree and centralization is fastest. In
the false-alarm case, the median rule fully filters the single
inaccurate report.

Three declared nonnegative scalarizations therefore select three
different winners: emergency-command weights select the centralized
architecture, adaptive-legitimacy weights select the polycentric
architecture, and filtering-and-stability weights select the
majoritarian architecture. The adaptive speed result is explicitly
feedback-dependent: with a nine-step feedback delay polycentric
correction ties the majoritarian case at thirteen updates, from ten
steps onward it is slower, and without per-group outcome feedback it
takes thirty-five updates. The construction demonstrates conditional
non-entailments, not the superiority of one political form.

\section{4. Institutions: continuity, legitimacy, and
correctability}\label{institutions-continuity-legitimacy-and-correctability}

\subsection{4.1 What institutions do}\label{what-institutions-do}

Institutions preserve coordination beyond the memory, intention, or
lifetime of any one participant. They stabilize roles, rules, records,
rights, obligations, resources, procedures, sanctions, and expectations.
A court, school, hospital, firm, legislature, union, monastery,
scientific society, or regulatory agency persists because some pattern
of authority and practice survives turnover in persons.

Institutional continuity is valuable when it carries forward capacities
that remain worth preserving. It is dangerous when continuity shields
the institution from evidence that its purposes, procedures, or effects
have changed.

\textbf{{[}Definition{]} Institutional continuity} is the persistence of
a role-and-rule organization through personnel and environmental change.
It is not by itself legitimacy, functional success, or fidelity to
purpose.

\subsection{4.2 Constitutive purpose and multiple
publics}\label{constitutive-purpose-and-multiple-publics}

An institution may have several purposes at once. A university educates,
certifies, researches, employs, allocates prestige, and reproduces
social networks. A hospital treats illness, trains practitioners,
manages risk, employs workers, bills payers, and satisfies regulation. A
platform connects users, sells attention, moderates content, and
generates data.

These purposes can conflict. There may be no single objective function
that faithfully represents them. Institutional analysis should therefore
distinguish:

\begin{itemize}
\tightlist
\item
  publicly stated purpose;
\item
  legal mandate;
\item
  operational incentives;
\item
  revenue model;
\item
  professional norms;
\item
  effects on affected groups;
\item
  and informal survival imperatives.
\end{itemize}

When an institution claims one purpose while its incentive structure
reliably produces another, the mismatch is not explained merely by
individual hypocrisy. It can be a property of the organization.

\subsection{4.3 Legitimacy as a layered
judgment}\label{legitimacy-as-a-layered-judgment}

Four layers should be distinguished:

\begin{enumerate}
\def\labelenumi{\arabic{enumi}.}
\tightlist
\item
  \textbf{Legal validity:} whether authority is exercised under
  recognized law.
\item
  \textbf{Sociological acceptance:} whether people believe the authority
  is proper or comply voluntarily.
\item
  \textbf{Procedural legitimacy:} whether decisions are made through
  fair, intelligible, participatory, and reviewable procedures.
\item
  \textbf{Normative legitimacy:} whether the institution's exercise of
  power can be justified to those subject to it, including those who
  lose.
\end{enumerate}

These layers can support one another without collapsing into identity. A
legally valid decision can be substantively unjust. A popular
institution can deny standing to a minority. A procedurally careful
system can operate within an unjust distribution of property or power.
An institution can perform well while preventing affected people from
contesting how ``performance'' is defined.

The provisional CU-S5 and CU-Inst-0 principles therefore treat
legitimacy as dependent on agency, legibility, contestability, burden
distribution, and repair---not on stability alone.

\subsection{4.4 Correctability}\label{correctability}

\textbf{{[}Definition{]} Institutional correctability} is the capacity
for relevant evidence, affected-person testimony, reasoned objection,
and observed consequences to produce review and, where warranted, change
in records, decisions, rules, roles, or remedies.

Correctability requires more than a complaint form. It may require:

\begin{itemize}
\tightlist
\item
  independent evidence channels;
\item
  protection against retaliation;
\item
  preservation of records;
\item
  transparent reasons;
\item
  review by actors not invested in the original decision;
\item
  authority to reverse or modify outcomes;
\item
  restitution or compensation;
\item
  and public learning from recurrent error.
\end{itemize}

An appeal process that cannot alter anything is expressive rather than
corrective. A transparency system that discloses information only after
irreversible harm may be too late. A whistleblowing policy without
protection changes the cost of truth-telling without changing the
underlying power relation.

Nominal channels do not establish practical correctability. Institutions
often control the categories, records, accessibility data, success
criteria, and recurrence measures by which their own review systems are
judged. Evidence of correctability should therefore include
disaggregated access and attrition, retaliation risk, timeliness,
reviewer independence, reversal and record-correction rates, remedy and
restitution, recurrence, and testimony from affected parties. No
institution should be the sole certifier of its own correctability;
external or jointly governed verification is required when the stakes
are substantial.

\subsubsection{4.4.1 Correctability, metrics, and the target
condition}\label{correctability-metrics-and-the-target-condition}

At institutional scale, correctability is lost when measures,
narratives, or review systems become insulated from the effects they are
meant to track. Metric absolutism, silent optimization, captured
records, and self-certified correction can improve the institutional
reading while the underlying social condition worsens. The diagnosis
requires evidence of omitted effects, asymmetric standards, blocked
challenge, or failure to repair. It is not inferred from disagreement,
poor outcomes, or institutional complexity alone.

\subsection{4.5 Role coherence}\label{role-coherence}

Institutional roles become unstable or abusive when authority,
information, capability, responsibility, and accountability are
misaligned.

Examples include:

\begin{itemize}
\tightlist
\item
  workers held responsible for outcomes while denied resources or
  discretion;
\item
  managers empowered to impose risk while insulated from its costs;
\item
  technical experts asked to certify decisions whose objectives were set
  elsewhere;
\item
  automated systems making consequential classifications without an
  accountable human role;
\item
  and elected bodies formally responsible for problems governed by
  actors outside their jurisdiction.
\end{itemize}

CU-ICF-2 names \textbf{responsibility without authority}. CU-ICF-3 names
\textbf{authority without accountability}. CU-Inst-6 requires their
realignment.

\subsection{4.6 Epistemic legibility and its
limits}\label{epistemic-legibility-and-its-limits}

Those affected by an institution should be able to understand the rules,
evidence, categories, objectives, and decision pathways that materially
govern them. This is epistemic legibility.

Legibility is not total publicity. Privacy, security, confidentiality,
and protection of vulnerable persons can justify limits. Nor is every
system fully explainable in simple terms. The requirement is
proportional: people need enough information to know how decisions are
made, challenge material errors, identify conflicts of interest, and
understand available remedies.

Legibility must also run upward. Frontline experience and local
knowledge need channels into institutional decision. A system that makes
populations legible to administrators while making administrators opaque
to populations creates one-way visibility.

\subsection{4.7 Scale alignment and
subsidiarity}\label{scale-alignment-and-subsidiarity}

A decision should be located at a scale capable of perceiving and
addressing its consequences without unnecessarily erasing local
knowledge or agency.

This is not a formula that always favors decentralization. Local
decisions can impose regional or planetary externalities. Central
decisions can miss local variation. Scale alignment asks:

\begin{itemize}
\tightlist
\item
  At what scale does the problem arise?
\item
  At what scale are consequences borne?
\item
  Where is relevant knowledge located?
\item
  Which actors can coordinate a response?
\item
  Who can contest error?
\item
  Which externalities cross the boundary?
\end{itemize}

Polycentric governance is attractive when multiple centers can act,
learn, and coordinate without one center monopolizing authority. But
polycentricity requires mechanisms for cross-scale conflict, resource
inequality, common standards, and accountability. Otherwise it can
become fragmentation with an attractive name.

\subsection{4.8 Institutional memory}\label{institutional-memory}

Records, precedents, routines, professional norms, and physical
infrastructure let institutions learn beyond individual tenure. They can
also preserve outdated categories and unacknowledged harm.

Institutional memory should therefore be:

\begin{itemize}
\tightlist
\item
  durable enough to prevent repetitive amnesia;
\item
  contestable enough to correct false records;
\item
  contextual enough to preserve why a decision was made;
\item
  accessible to those affected;
\item
  and governed by privacy and expiration where indefinite retention
  becomes harmful.
\end{itemize}

Memory is not inherently conservative or progressive. It is an
architecture of temporal power.

\subsection{4.9 Participation and
tokenism}\label{participation-and-tokenism}

Participation matters when it changes the information, alternatives,
criteria, or authority of a decision. Consultation can become ritualized
extraction when institutions ask affected people to contribute
experience but retain unilateral control over interpretation and
outcome.

Substantive participation may require:

\begin{itemize}
\tightlist
\item
  agenda-setting authority;
\item
  compensation for time and expertise;
\item
  accessible formats;
\item
  language support;
\item
  representation of internally diverse views;
\item
  disclosure of what is and is not negotiable;
\item
  and a public account of how contributions altered the decision.
\end{itemize}

Affected communities are not epistemically infallible and rarely speak
with one voice. Their standing arises because they possess relevant
knowledge and bear consequences, not because identity confers automatic
correctness.

\subsection{4.10 Institutional failure
diagnostics}\label{institutional-failure-diagnostics}

Five provisional diagnostics are especially useful.

\textbf{CU-ICF-1 --- Metric Absolutism.} A proxy replaces the plural
purposes and constraints it was introduced to serve.

\textbf{CU-ICF-2 --- Responsibility Without Authority.} Agents are held
answerable for outcomes they lack the information, resources,
discretion, or power to change.

\textbf{CU-ICF-3 --- Authority Without Accountability.} Decision power
is concentrated without proportionate transparency, answerability,
correction, or repair.

\textbf{CU-ICF-4 --- Normalization Pressure.} Functional diversity and
local knowledge are erased to simplify measurement, management, or
enforcement.

\textbf{CU-ICF-5 --- Silent Optimization.} Material tradeoffs are
imposed through technical or administrative processes without making the
objective, affected parties, and alternatives publicly contestable.

These diagnostics identify structural risks. They do not prove intent or
moral culpability in every case.

\section{5. Power, material security, and
extraction}\label{power-material-security-and-extraction}

\subsection{5.1 Why material conditions are not external to social
coherence}\label{why-material-conditions-are-not-external-to-social-coherence}

Coordination requires bodies, time, shelter, care, energy,
communication, mobility, and freedom from constant emergency. Material
insecurity narrows practical options and can make nominal rights
unusable. A person may possess a formal right to appeal while lacking
transportation, childcare, legal assistance, documentation, or time away
from work.

The ILO's 2024--2026 social-protection report and the United Nations'
2025 social report document large protection gaps and emphasize the
relationship among insecurity, inequality, capability, and public trust
(ILO, 2024; United Nations Department of Economic and Social Affairs,
2025). These reports do not establish one global social-coherence
variable. They support the narrower claim that material floors are part
of the infrastructure through which people can participate, recover from
shocks, and exercise agency.

\subsection{5.2 Capabilities rather than resources
alone}\label{capabilities-rather-than-resources-alone}

Resources matter through what people are actually able to do and be.
Equal formal allocation can produce unequal capability when people face
different bodily, social, geographic, or caregiving conditions. The
capability approach developed by Sen and Nussbaum is therefore
compatible with this paper's insistence on affected-person scope and
plural dimensions (Sen, 1999; Nussbaum, 2011).

A social profile should distinguish:

\begin{itemize}
\tightlist
\item
  income and wealth;
\item
  access to health, education, housing, and transport;
\item
  time autonomy;
\item
  security against catastrophic loss;
\item
  political and associational capability;
\item
  and the practical ability to refuse harmful arrangements.
\end{itemize}

\subsection{5.3 Extraction}\label{extraction}

\textbf{{[}Definition{]} Extraction} occurs when an actor or system
captures benefit by transferring substantial nonconsensual, hidden,
uncompensated, or effectively unrepairable costs to other persons,
communities, future generations, or ecosystems.

Extraction can occur through:

\begin{itemize}
\tightlist
\item
  low wages combined with high dependency;
\item
  pollution externalized onto communities;
\item
  unpaid or underrecognized care work;
\item
  data and attention capture;
\item
  predatory debt;
\item
  legal asymmetries;
\item
  monopoly control;
\item
  depletion of common resources;
\item
  or institutional arrangements in which those who benefit can exit
  before costs arrive.
\end{itemize}

Not every unequal exchange is extraction. The diagnosis depends on
consent, alternatives, power, risk, information, repair, and the
distribution of gains and losses.

\subsection{5.4 Local efficiency and global
burden}\label{local-efficiency-and-global-burden}

An organization can become more efficient by exporting costs. A platform
can increase engagement while degrading attention. A firm can lower
prices while shifting environmental or labor costs outside its accounts.
A state can stabilize domestic conditions through burdens imposed
elsewhere.

Accordingly:

\[\adjustbox{max width=\linewidth}{$\displaystyle 
\text{local performance improvement}
\not\Rightarrow
\text{global social improvement}.
$}\]

This is the social version of the local--global discipline inherited
from Biology and Psychology. It does not claim a universal conservation
law. It requires the analyst to trace where costs go.

\subsection{5.5 Inequality, difference, and
domination}\label{inequality-difference-and-domination}

Inequality is not one phenomenon. Differences in income, wealth, status,
legal standing, health, exposure to violence, political influence, time,
and epistemic credibility have different mechanisms and consequences.

The key question is when unequal positions enable one party to set the
terms under which another must live while remaining insulated from
reciprocal influence. This is domination.

\textbf{{[}Definition{]} Domination} is durable exposure to another's
capacity for arbitrary interference or rule setting without adequate
reciprocal constraint, contestation, or exit.

A relationship can be benevolent in intention and dominating in
structure. Conversely, not every hierarchy is arbitrary: emergency
command, expertise, guardianship, and delegated authority can be
legitimate under bounded conditions. The burden is to explain scope,
accountability, review, and termination.

\subsection{5.6 Economic coordination}\label{economic-coordination}

Markets are powerful mechanisms for coordinating exchange under certain
institutional conditions. They communicate some information through
prices and allow decentralized adaptation. They do not generate their
own moral boundaries, represent every value, or ensure that willingness
to pay tracks need, contribution, or justice.

This paper treats markets as bounded coordination mechanisms among
others. Their assessment includes:

\begin{itemize}
\tightlist
\item
  market structure and concentration;
\item
  bargaining power;
\item
  information asymmetry;
\item
  externalities;
\item
  public goods;
\item
  social insurance;
\item
  labor conditions;
\item
  ecological limits;
\item
  and the distribution of ownership and risk.
\end{itemize}

Growth, productivity, revenue, and capital formation are instrumental
achievements. Their social value depends on what they enable, whom they
benefit, what they displace, and whether repair costs are counted.

\subsection{5.7 Labor and agency}\label{labor-and-agency}

Labor involves persons whose bodies, time, attention, skill,
relationships, and future options have non-instrumental significance.
Employment can create capability, identity, belonging, and material
security. It can also generate injury, precarity, surveillance, and
dependency.

A job cannot be assessed solely through wage or output. Relevant
dimensions include:

\begin{itemize}
\tightlist
\item
  control over time and pace;
\item
  bodily and psychological safety;
\item
  voice;
\item
  predictability;
\item
  learning;
\item
  recognition;
\item
  collective bargaining;
\item
  care obligations;
\item
  and ability to leave without catastrophic loss.
\end{itemize}

\subsection{5.8 Attention as a social
resource}\label{attention-as-a-social-resource}

Attention is limited and developmentally significant. Digital systems
can profit by selecting, ranking, and repeating stimuli that capture
attention, shape desire, and reorganize public salience.

The economic value of attention does not make every attention-maximizing
design legitimate. CU-EC-10 treats autonomous attention as a protected
scarce capacity. Governance questions include:

\begin{itemize}
\tightlist
\item
  what is optimized;
\item
  how content is selected;
\item
  whether users can understand and alter the process;
\item
  effects on children and vulnerable populations;
\item
  and whether business incentives reward epistemic or psychological
  harm.
\end{itemize}

\subsection{5.9 Repair obligations}\label{repair-obligations}

Where an activity generates foreseeable harm, repair costs should not be
treated as external to the activity merely because another institution
pays them later. Those who capture substantial benefits may bear
obligations proportional to causal contribution, capacity, knowledge,
and control.

This is a normative proposal, not a physical conservation law. It
requires legal and ethical argument about attribution, foreseeability,
collective responsibility, and fair procedure.

\section{6. Information integrity, culture, and public
meaning}\label{information-integrity-culture-and-public-meaning}

\subsection{6.1 Information integrity}\label{information-integrity}

\textbf{{[}Definition{]} Information integrity} is the capacity of a
social system to produce, preserve, contest, transmit, and correct
claims with sufficient provenance and independence for collective
decisions to remain answerable to reality.

Information integrity is not universal agreement. It includes:

\begin{itemize}
\tightlist
\item
  access to evidence;
\item
  source provenance;
\item
  professional and methodological standards;
\item
  independent institutions;
\item
  adversarial checking;
\item
  correction mechanisms;
\item
  public records;
\item
  protection for investigative work;
\item
  and capacity to distinguish uncertainty from fabrication.
\end{itemize}

A system can contain abundant information while lacking integrity if
visibility is controlled by incentives unrelated to accuracy.

\subsection{6.2 Shared reality without one authorized
narrative}\label{shared-reality-without-one-authorized-narrative}

Collective action often requires agreement about practical facts:
whether a bridge is safe, an election result was counted, a pathogen is
spreading, or a budget contains sufficient funds. It does not require
agreement about every moral, metaphysical, cultural, or political
interpretation.

The aim is therefore not one total worldview. It is a layered public
epistemology in which:

\begin{itemize}
\tightlist
\item
  factual claims remain corrigible;
\item
  uncertainty is representable;
\item
  disagreements can be localized;
\item
  evidence channels remain plural;
\item
  and no actor acquires unilateral authority to define all legitimate
  knowledge.
\end{itemize}

\subsection{6.3 Misinformation, disinformation, and
error}\label{misinformation-disinformation-and-error}

The following should be distinguished:

\begin{itemize}
\tightlist
\item
  \textbf{misinformation:} false or misleading content shared without
  requiring deceptive intent;
\item
  \textbf{disinformation:} false or misleading content produced or
  distributed strategically;
\item
  \textbf{malinformation:} genuine information used out of context or in
  ways designed to harm;
\item
  \textbf{ordinary error:} inaccurate belief arising without strategic
  manipulation;
\item
  and \textbf{contested interpretation:} disagreement not settled merely
  by attaching a truth label.
\end{itemize}

Interventions appropriate to one category may fail in another.
Correction, inoculation, source labeling, friction, moderation,
transparency, public-interest media, and changes to platform incentives
act on different mechanisms.

\subsection{6.4 Networks and
amplification}\label{networks-and-amplification}

Social networks can improve access to diverse knowledge and mutual aid.
They can also amplify correlated error and motivational bias. In
controlled experiments, transmission through social networks has been
shown to amplify biased judgments under some conditions, while more
representative resampling reduced that amplification without eliminating
the benefits of social information (Hardy et al., 2023).

The lesson is not that networks are inherently distorting. It is that
sampling architecture matters. The visible stream is not a neutral
window onto public belief.

\subsection{6.5 Culture as
infrastructure}\label{culture-as-infrastructure}

Culture supplies symbols, narratives, categories, rituals, expectations,
identities, and practices through which collective life becomes
intelligible. It reduces some coordination costs by making roles and
obligations predictable. It also shapes desire, shame, aspiration,
belonging, and what kinds of life appear possible.

\textbf{{[}Bridge principle{]} Culture is meaning and coordination
infrastructure.} This does not make culture one field or one memory
store. Cultures are internally contested, layered, hybrid, and
distributed across institutions and media.

\subsection{6.6 Cultural capacity and cultural
harm}\label{cultural-capacity-and-cultural-harm}

Cultural practices can create capacities for:

\begin{itemize}
\tightlist
\item
  mourning;
\item
  care;
\item
  mutual aid;
\item
  stewardship;
\item
  intergenerational memory;
\item
  moral imagination;
\item
  skill;
\item
  celebration;
\item
  and non-instrumental belonging.
\end{itemize}

They can also normalize:

\begin{itemize}
\tightlist
\item
  domination;
\item
  caste or racial hierarchy;
\item
  exhaustion;
\item
  consumption;
\item
  gendered self-erasure;
\item
  punitive purity;
\item
  and sacrifice of outsiders.
\end{itemize}

A tradition's longevity does not settle its legitimacy. A practice's
sacred meaning does not remove its effects from evaluation. Conversely,
the presence of historical harm does not imply that every inherited
practice is reducible to domination.

\subsection{6.7 Narrative capture}\label{narrative-capture}

\textbf{{[}Definition{]} Narrative capture} occurs when dominant
meanings make alternatives unintelligible, shameful, professionally
dangerous, or socially catastrophic rather than openly contestable.

Capture can operate through state propaganda, corporate branding,
ideological movements, professional jargon, or intimate group pressure.
It is not established merely because a narrative is popular or
emotionally powerful. The diagnostic question is whether relevant
counterevidence and alternative interpretations can enter without
automatic loss of standing.

\subsection{6.8 Meaning renewal}\label{meaning-renewal}

People and communities need ways to absorb loss, role transition,
migration, failure, historical disclosure, and technological change.
Meaning that can survive only by denying reality is brittle. Meaning
that cannot survive any revision is coercive. Meaning that is
continually deconstructed without practices of renewal leaves people
without orientation.

CU-CL-8 therefore proposes that humane cultures should permit meaning to
be regenerated after loss and change. This is a normative commitment,
not a measurable law of cultural evolution.

\subsection{6.9 Refusal and exit}\label{refusal-and-exit}

A cultural narrative is safer when participants can reinterpret roles,
dissent, leave, or return without losing basic civil or material
standing. Exit alone is insufficient when costs are catastrophic or when
all alternatives are controlled by the same system. Voice alone is
insufficient when it cannot alter decisions.

Hirschman's distinction among exit, voice, and loyalty remains useful
here (Hirschman, 1970). The further point is that their availability and
cost are distributed unequally.

\section{7. Polarization, conflict, and democratic
correction}\label{polarization-conflict-and-democratic-correction}

\subsection{7.1 Polarization is not one
variable}\label{polarization-is-not-one-variable}

The term \emph{polarization} is used for several different phenomena:

\begin{itemize}
\tightlist
\item
  ideological distance on policy;
\item
  affective hostility toward an outgroup;
\item
  partisan sorting across identities;
\item
  elite polarization;
\item
  perceived extremity of opponents;
\item
  segregation of information networks;
\item
  moralization of political identity;
\item
  and willingness to violate democratic constraints for partisan ends.
\end{itemize}

These dimensions can move separately. An intervention that warms
feelings toward opponents may not increase support for democratic rules.
Voelkel and colleagues found that interventions reducing affective
polarization did not necessarily reduce anti-democratic attitudes,
support for partisan violence, or willingness to privilege partisan ends
over democratic means (Voelkel et al., 2023). The result is a direct
warning against treating one social indicator as a sufficient proxy for
institutional health.

\subsection{7.2 Conflict can be
informative}\label{conflict-can-be-informative}

Political conflict may reflect manipulation, identity threat, and
institutional failure. It may also reflect genuine disagreement over
distribution, rights, recognition, history, and risk. A society in which
marginalized groups cease protesting may have resolved injustice,
exhausted resistance, or increased repression. Surface calm does not
identify which.

The framework therefore asks whether conflict is:

\begin{itemize}
\tightlist
\item
  representable without dehumanization;
\item
  connected to evidence and affected interests;
\item
  governed by legitimate procedures;
\item
  constrained against violence and retaliation;
\item
  and capable of producing policy revision or negotiated coexistence.
\end{itemize}

\subsection{7.3 Affective polarization and democratic
norms}\label{affective-polarization-and-democratic-norms}

Partisan animosity can weaken cooperation and support for democratic
restraint under some conditions, but causal pathways vary. Reviews
identify interventions at the levels of beliefs, relationships, elite
behavior, and institutions, with substantial uncertainty concerning
durability, scaling, and backfire (Hartman et al., 2022).

Accordingly, ``depolarization'' should not be the sole goal. Some
interventions may reduce visible disagreement by suppressing political
participation or minimizing real injustice. The protected outcomes are
nonviolence, equal standing, truthful public reasoning, fair
competition, and the ability to revise power---not emotional warmth by
itself.

\subsection{7.4 Misperception and self-fulfilling
polarization}\label{misperception-and-self-fulfilling-polarization}

Groups often overestimate the extremity, hostility, or anti-democratic
attitudes of opponents. Such beliefs can become self-fulfilling when
actors pre-emptively abandon restraint. Correcting misperceptions can
help in some contexts, but it cannot substitute for institutional reform
where incentives genuinely reward extremity.

The analyst should therefore distinguish:

\begin{itemize}
\tightlist
\item
  false beliefs about opponents;
\item
  accurate beliefs about elite strategy;
\item
  symbolic conflict;
\item
  material conflict;
\item
  and institutional structures that convert narrow electoral advantage
  into concentrated power.
\end{itemize}

\subsection{7.5 Democratic institutions as correction
systems}\label{democratic-institutions-as-correction-systems}

Democracy is not only a method for aggregating preference. At its
strongest, it distributes the capacity to replace rulers, contest
policy, expose information, organize opposition, protect civil
liberties, and revise institutions without civil war.

This makes democratic institutions one family of social correction
architectures. Their performance depends on:

\begin{itemize}
\tightlist
\item
  meaningful competition;
\item
  freedom of expression and association;
\item
  independent media and civil society;
\item
  equal legal standing;
\item
  checks on executive power;
\item
  credible administration of elections;
\item
  and the practical ability of affected groups to participate.
\end{itemize}

The V-Dem 2026 report documents a broad global pattern of
autocratization and declining liberal-democratic conditions through
2025. Its indices and expert-coded data provide one important evidence
stream, not a complete or politically neutral measure of every society
(Nord et al., 2026). This paper uses this evidence to motivate analysis
of correction capacity, not to collapse democracy into one
social-coherence score.

\subsection{7.6 Autocratization as erosion of
correctability}\label{autocratization-as-erosion-of-correctability}

Autocratization often proceeds through incremental changes rather than
immediate abolition of elections. Mechanisms can include weakening
legislative and judicial constraints, restricting media and civil
society, altering electoral administration, politicizing enforcement,
concentrating executive authority, and raising the cost of opposition.

From this perspective, the common structure is degradation of
independent correction channels. This does not mean every institutional
dispute is autocratization. The diagnosis requires comparative evidence,
legal and historical context, and attention to actual changes in power.

\subsection{7.7 Repression and preference
falsification}\label{repression-and-preference-falsification}

When dissent is costly, public agreement ceases to provide reliable
information about private belief. People may conform outwardly while
privately rejecting the regime, producing sudden cascades when
expectations about others change (Kuran, 1995).

This creates an epistemic problem for authorities as well as subjects. A
system that punishes bad news can lose access to the information
required for its own survival. Repression can produce apparent stability
while increasing uncertainty about underlying support.

\subsection{7.8 Democratic resilience without
romanticism}\label{democratic-resilience-without-romanticism}

Democratic systems can be slow, unequal, captured, exclusionary, and
vulnerable to misinformation. Protected contestation is not sufficient
for good outcomes. It must be joined to material capacity, public
administration, anti-corruption safeguards, rights, civic organization,
and institutions able to implement decisions.

The relevant contrast is therefore not ideal democracy versus defective
alternatives. It is among actual arrangements with different capacities
for coordination, representation, correction, and restraint.

\subsection{7.9 Violence and the failure of ordinary
correction}\label{violence-and-the-failure-of-ordinary-correction}

Political violence becomes more likely through multiple pathways: state
repression, armed organization, security dilemmas, impunity,
dehumanization, territorial control, economic shocks, and breakdown of
credible institutions. No single coherence mechanism explains these
outcomes.

The narrower claim is that violence becomes structurally more likely
when actors believe that ordinary channels cannot protect them, correct
error, or enforce agreements. Restoring those channels may be necessary
but is not always sufficient.

\section{8. Meso-scale social
infrastructure}\label{meso-scale-social-infrastructure}

\subsection{8.1 Why the middle scale
matters}\label{why-the-middle-scale-matters}

Between the individual and the state lie the structures through which
most social regulation occurs: households, friendships, neighborhoods,
congregations, schools, workplaces, unions, professional associations,
cooperatives, mutual-aid networks, local media, clubs, and civic
organizations.

These structures provide:

\begin{itemize}
\tightlist
\item
  belonging and recognition;
\item
  informal care;
\item
  practical knowledge;
\item
  identity and ritual;
\item
  apprenticeship;
\item
  conflict mediation;
\item
  early warning of distress;
\item
  and channels between persons and larger institutions.
\end{itemize}

The meso scale is a major site of both social erosion and repair, but
dense community is not automatically beneficial.

\subsection{8.2 Bonding, bridging, and linking
relations}\label{bonding-bridging-and-linking-relations}

Social-capital research often distinguishes:

\begin{itemize}
\tightlist
\item
  \textbf{bonding ties} within relatively similar groups;
\item
  \textbf{bridging ties} across social differences; and
\item
  \textbf{linking ties} connecting people to institutions and unequal
  positions of authority.
\end{itemize}

Each can support or obstruct collective life. Bonding can provide deep
care while intensifying exclusion. Bridging can expand trust while
remaining superficial. Linking can provide access to institutions while
reproducing patronage.

The quality, power, and purpose of relations matter more than tie count
alone.

\subsection{8.3 Social connection and
health}\label{social-connection-and-health}

The WHO Commission on Social Connection's 2025 report treats loneliness
and social isolation as widespread public-health and social concerns and
calls for action across policy, community, organizational, and
individual levels (WHO Commission on Social Connection, 2025). The
general lesson imported here is that social connection is materially and
psychologically consequential.

It does not infer that every person needs the same amount or form of
contact, that solitude is pathological, or that institutional belonging
is always safe. Connection should be evaluated by quality,
voluntariness, accessibility, reciprocity, and fit.

\subsection{8.4 Families and intimate
networks}\label{families-and-intimate-networks}

Families can provide care, memory, material support, and regulation.
They can also distribute labor and authority unequally, hide abuse, and
impose identities. Policy that assumes family support without examining
its availability can privatize public obligations and intensify
dependency.

The relevant questions include:

\begin{itemize}
\tightlist
\item
  Who provides care?
\item
  Who controls resources?
\item
  Who can leave?
\item
  Which relationships are legally recognized?
\item
  What support exists when family is absent or unsafe?
\end{itemize}

\subsection{8.5 Congregations and communities of
practice}\label{congregations-and-communities-of-practice}

Religious and contemplative communities can sustain ritual, moral
formation, service, grief practices, intergenerational contact, and
non-market belonging. They can also enforce conformity, conceal harm, or
make spiritual standing dependent on obedience.

Their social role cannot be evaluated from doctrine alone. Governance,
leadership accountability, gender and status relations, handling of
abuse, and capacity for dissent are equally relevant.

\subsection{8.6 Workplaces and
organizations}\label{workplaces-and-organizations}

Workplaces organize a large share of adult time, identity, learning, and
social contact. They can provide stable relations and shared purpose or
intensify precarity, surveillance, competition, and burnout.

An organization's internal coherence is not established by productivity
or mission alignment. It depends also on whether people can report
error, whether roles are realistic, whether burdens are repairable, and
whether authority is accountable.

\subsection{8.7 Movements and
counter-institutions}\label{movements-and-counter-institutions}

Social movements can reveal harms that established institutions do not
perceive or prefer not to acknowledge. They can create new identities,
knowledge, solidarity, and political possibility. They can also become
purity-driven, centralized, factional, or detached from the people they
claim to represent.

A movement's oppositional position does not exempt it from the same
questions of power, evidence, accountability, and repair.

\subsection{8.8 Meso-scale erosion}\label{meso-scale-erosion}

Meso institutions can weaken through residential instability, long work
hours, declining public space, economic precarity, institutional
distrust, technological substitution, demographic change, and deliberate
repression. The causal mix differs by location.

Putnam's account of declining civic participation in the United States
remains influential, but its measures and historical interpretation are
not universal templates (Putnam, 2000). Contemporary social connection
must include digital, diasporic, mutual-aid, disability, and
identity-based forms that older civic indicators may miss.

\subsection{8.9 Rebuilding without
nostalgia}\label{rebuilding-without-nostalgia}

Rebuilding meso-scale social infrastructure does not mean restoring one
idealized community form. Past communities often depended on exclusion,
unpaid labor, enforced conformity, and restricted mobility.

Reconstruction should instead ask which functions need support:

\begin{itemize}
\tightlist
\item
  dependable care;
\item
  accessible gathering space;
\item
  local journalism;
\item
  democratic association;
\item
  intergenerational contact;
\item
  shared skill and ritual;
\item
  conflict mediation;
\item
  and links to formal institutions.
\end{itemize}

Different communities can realize these functions differently.

\section{9. Crisis, cascades, resilience, and
thresholds}\label{crisis-cascades-resilience-and-thresholds}

\subsection{9.1 From one crisis to interacting
processes}\label{from-one-crisis-to-interacting-processes}

The phrase \emph{coherence crisis} can orient attention to
interdependence, but it becomes misleading if treated as a measured
diagnosis. Current societies face heterogeneous and uneven pressures:

\begin{itemize}
\tightlist
\item
  material insecurity and unequal exposure to shocks;
\item
  declining trust in some institutions and groups;
\item
  democratic backsliding in many countries alongside democratization in
  others;
\item
  social isolation and loneliness;
\item
  ecological disruption;
\item
  migration and displacement;
\item
  armed conflict;
\item
  information disorder;
\item
  concentrated technological power;
\item
  and weakened capacity for long-horizon action.
\end{itemize}

The United Nations' 2025 social report connects insecurity, inequality,
distrust, and pressure on solidarity while arguing for equity, economic
security, and solidarity as policy priorities (United Nations Department
of Economic and Social Affairs, 2025). The report is a global policy
synthesis, not proof that one mechanism explains every trend.

\subsection{9.2 Cascades}\label{cascades}

A cascade occurs when a change in one part of a network alters
conditions for others in a way that propagates. Examples include bank
runs, protest diffusion, supply-chain interruption, rumor spread, policy
imitation, violence escalation, and failures in tightly coupled
infrastructure.

Cascade claims require:

\begin{itemize}
\tightlist
\item
  a network or dependency structure;
\item
  a transmission mechanism;
\item
  thresholds or response functions;
\item
  a time scale;
\item
  and evidence separating propagation from common external cause.
\end{itemize}

Two events moving together do not establish a cascade.

\subsection{9.3 Resilience}\label{resilience}

\textbf{{[}Definition{]} Social resilience} is the capacity of a
declared social unit to prepare for, withstand, adapt to, and recover
from a specified class of shocks while preserving protected conditions
and the standing of affected groups.

This definition follows the boundary--shock--horizon discipline
emphasized in recent applied resilience frameworks (World Bank, 2024).
It rejects the idea that resilience is a generic trait.

A regime can appear resilient by transferring harm to marginalized
groups. A household can ``cope'' by withdrawing a child from school. A
city can restore aggregate output while displaced residents cannot
return. Resilience must therefore include distribution and
transformation, not only recovery of system averages.

\subsection{9.4 Robustness, adaptability, and
transformation}\label{robustness-adaptability-and-transformation}

These terms should be separated.

\begin{itemize}
\tightlist
\item
  \textbf{Robustness:} maintaining function under disturbance without
  major structural change.
\item
  \textbf{Adaptability:} modifying behavior, allocation, or rules while
  preserving the basic institution.
\item
  \textbf{Transformability:} changing the institution or system when the
  previous structure can no longer meet protected constraints.
\end{itemize}

A system can be robust against ordinary variation and unable to
transform when its environment changes. Conversely, continual
transformation can destroy reliability.

\subsection{9.5 Thresholds and regime
shifts}\label{thresholds-and-regime-shifts}

Some social processes exhibit nonlinear thresholds. Participation
cascades, financial contagion, conflict escalation, and institutional
breakdown can change rapidly when expectations or network conditions
cross a boundary.

This paper does not infer a universal social phase transition. The
existence, location, and meaning of a threshold depend on the model and
observation scale. ``Collapse'' must specify what has failed: government
continuity, public services, legitimacy, monetary coordination, peaceful
succession, ecological support, or some other condition.

\subsection{9.6 Early-warning signals}\label{early-warning-signals}

The critical-transition literature identifies possible warning patterns
such as rising variance, autocorrelation, skewness, or recovery time in
systems approaching certain bifurcations (Scheffer et al., 2009).
Transferring these indicators to social systems requires caution.

Social actors learn, anticipate measurement, change strategy, and alter
the system being observed. Time series are often short, nonstationary,
politically manipulated, and affected by multiple shocks. A warning
indicator should therefore report:

\begin{itemize}
\tightlist
\item
  the assumed mechanism;
\item
  baseline and comparison cases;
\item
  false positives and false negatives;
\item
  data revisions;
\item
  and whether the signal would still be meaningful after actors respond
  to it.
\end{itemize}

Rising variance is not a universal sign that society is near collapse.

\subsection{9.7 Hysteresis and path
dependence}\label{hysteresis-and-path-dependence}

The route into a crisis may differ from the route out. Trust lost
through betrayal may not return when performance improves. Displaced
communities do not reappear when prices fall. Emergency institutions
create interests and capacities that outlast the emergency. Violence
changes identities and expectations.

Hysteresis is therefore a useful hypothesis when a mechanism preserves
the effects of history. It should not be invoked merely because recovery
is difficult.

\subsection{9.8 Trauma across social
scales}\label{trauma-across-social-scales}

Collective violence, dispossession, racism, forced migration,
institutional betrayal, and economic deprivation can shape psychological
regulation, family relations, community trust, and political behavior.
But ``social trauma'' can become too broad if every conflict is
redescribed clinically.

Four restrictions follow:

\begin{enumerate}
\def\labelenumi{\arabic{enumi}.}
\tightlist
\item
  Individual diagnosis is not inferred from group membership.
\item
  Political disagreement is not treated as symptom.
\item
  Institutional and material mechanisms retain explanatory priority
  where appropriate.
\item
  Affected communities participate in defining injury and repair.
\end{enumerate}

Social conditions can create trauma exposure without the society itself
being a mind that literally has PTSD.

\subsection{9.9 Crisis response and
distribution}\label{crisis-response-and-distribution}

Emergency policy often reveals whose lives and work are treated as
expendable. Aggregate recovery indicators can conceal unequal mortality,
displacement, debt, educational loss, and long-term disability.

Crisis analysis should therefore track:

\begin{itemize}
\tightlist
\item
  exposure;
\item
  capacity to comply with protective measures;
\item
  access to relief;
\item
  decision-making voice;
\item
  long-run burden;
\item
  and who captures recovery benefits.
\end{itemize}

\subsection{9.10 No universal collapse
sequence}\label{no-universal-collapse-sequence}

Recurrent social-collapse sequences and competing attractor basins may
be useful prompts for case-specific modeling, but they are not
established social laws.

A regime may fail through fiscal exhaustion, war, succession crisis,
ecological shock, elite fragmentation, public-health catastrophe,
legitimacy collapse, infrastructure failure, or external conquest.
Similar surface outcomes can have different causal architectures.
Comparative work should search for recurrent mechanisms without forcing
every case into one trajectory.

\section{10. Social repair}\label{social-repair}

\subsection{10.1 Repair is not return}\label{repair-is-not-return}

Social repair is often imagined as restoration of a previous
equilibrium. That can be impossible or undesirable. The prior state may
have contained the conditions that produced the harm.

\textbf{{[}Definition{]} Social repair} is the process by which a
collective addresses consequential harm or institutional failure through
combinations of truth, safety, accountability, restitution, rule
revision, material reconstruction, restored standing, and renewed
capacity for coordinated life.

Repair may create a new arrangement rather than recover the old one.

\subsection{10.2 A repair sequence}\label{a-repair-sequence}

The following sequence is analytic, not universally chronological:

\begin{enumerate}
\def\labelenumi{\arabic{enumi}.}
\tightlist
\item
  \textbf{Cessation and protection:} stop ongoing harm and secure those
  at risk.
\item
  \textbf{Fact finding:} establish what happened, what remains
  uncertain, and which records require correction.
\item
  \textbf{Recognition:} acknowledge affected persons as knowers and
  bearers of standing.
\item
  \textbf{Responsibility:} identify causal contribution, authority,
  negligence, role, and capacity without reducing collective processes
  to one villain.
\item
  \textbf{Restitution or compensation:} restore resources, rights,
  records, opportunity, or material conditions where possible.
\item
  \textbf{Institutional change:} alter incentives, roles, rules,
  oversight, and evidence channels that enabled recurrence.
\item
  \textbf{Relational work:} where desired and safe, rebuild trust or
  workable coexistence.
\item
  \textbf{Memory and non-recurrence:} preserve enough history to prevent
  denial while avoiding permanent reduction of persons to the harm.
\end{enumerate}

Not every case requires every element. The sequence prevents apology
from standing in for material repair and prevents material settlement
from erasing truth.

\subsection{10.3 Justice, forgiveness, reconciliation, and
trust}\label{justice-forgiveness-reconciliation-and-trust}

These concepts are distinct.

\begin{itemize}
\tightlist
\item
  \textbf{Justice} concerns justified response to wrong, rights,
  burdens, and standing.
\item
  \textbf{Forgiveness} is a possible transformation in the injured
  person's relation to the wrong or wrongdoer.
\item
  \textbf{Reconciliation} is a renewed relation between parties or
  groups.
\item
  \textbf{Trust} is renewed vulnerability based on an expectation of
  trustworthiness.
\item
  \textbf{Coexistence} can be achieved without forgiveness,
  reconciliation, or trust.
\end{itemize}

No institution should require forgiveness as proof of repair. Renewed
trust must answer to changed conditions, not ceremonial pressure.

\subsection{10.4 Truth processes}\label{truth-processes}

Truth commissions, inquiries, audits, archives, investigative
journalism, testimony, and public history can expose patterns that
ordinary adjudication misses. Their effects vary with mandate, power,
implementation, political context, and relation to courts and material
reform. Reviews of truth-commission impact find mixed and
context-dependent results rather than one universal path from disclosure
to reconciliation (Bakiner, 2021).

Truth matters even when it does not immediately reconcile. At the same
time, truth processes can extract testimony, reopen injury, or produce
records without remedy. Design must include consent, support, security,
and control over use.

\subsection{10.5 Accountability}\label{accountability}

Accountability is not identical with punishment. It includes
answerability, exposure of reasons and evidence, consequences
proportionate to role and power, correction of records, removal from
authority where necessary, and obligations to repair.

Purely individual accountability can obscure structural causes. Purely
structural explanation can erase agency. Both levels are required where
both are causally relevant.

\subsection{10.6 Restitution and
redistribution}\label{restitution-and-redistribution}

Some harms are not repaired by recognition alone. Stolen land, unpaid
labor, discriminatory exclusion, destroyed infrastructure, environmental
contamination, and denied education have material effects.

Restitution may involve return, compensation, debt relief, service
guarantees, land or resource access, institutional investment, or
changes in ownership and control. The form must be argued case by case
with affected communities.

\subsection{10.7 Institutional learning}\label{institutional-learning}

An institution has not learned merely because it issues a report.
Learning requires changed behavior under relevant future conditions.

Evidence of institutional learning can include:

\begin{itemize}
\tightlist
\item
  changed decision rights;
\item
  revised incentives;
\item
  independent oversight;
\item
  safer reporting channels;
\item
  changed resource allocation;
\item
  measurable reduction in recurrence;
\item
  and public explanation of why the new structure should work.
\end{itemize}

\subsection{10.8 Repair and
irreversibility}\label{repair-and-irreversibility}

Some losses cannot be restored: death, lost years, ecological
extinction, destroyed archives, displaced communities, and developmental
injury. Repair remains meaningful but changes form. It may involve
mourning, memorialization, compensation, guarantees of non-recurrence,
and creation of capacities that did not previously exist.

The inability to undo harm increases the obligation to acknowledge
rather than deny it.

\subsection{10.9 Participation in repair}\label{participation-in-repair}

Those harmed should have substantive influence over how injury and
remedy are defined. This does not grant any one representative total
authority over a heterogeneous group. It requires structures that can
hold internal disagreement, unequal vulnerability, and changing needs.

\subsection{10.10 Repair as capacity}\label{repair-as-capacity}

A socially repairable institution can:

\begin{itemize}
\tightlist
\item
  detect mismatch;
\item
  keep evidence accessible;
\item
  protect reporters;
\item
  pause harmful action;
\item
  revise decisions;
\item
  compensate injury;
\item
  change the generating structure;
\item
  and restore standing without demanding silence.
\end{itemize}

Repair capacity is not the same as low conflict. Institutions with
visible appeals may record more complaints precisely because people
believe correction is possible.

\section{11. Technology, platforms, and AI-mediated social
systems}\label{technology-platforms-and-ai-mediated-social-systems}

\subsection{11.1 Technology is
institutional}\label{technology-is-institutional}

Technologies enter social life through ownership, business models,
interfaces, standards, law, labor, and administrative use. A platform is
not only an algorithm. It is an institution that defines participation,
visibility, moderation, identity, data access, and appeal.

Technological effects therefore cannot be inferred from code alone. They
depend on deployment, user adaptation, competition, culture, regulation,
and complementary institutions.

\subsection{11.2 Ranking and public
salience}\label{ranking-and-public-salience}

Feed, search, recommendation, and advertising systems determine which
items become visible and in what order. Ranking can affect attention and
sometimes attitudes, but effects vary by platform and intervention.

The 2020 Facebook and Instagram election studies found that substantial
feed changes altered exposure and engagement without producing
detectable changes in several measured political attitudes or
polarization outcomes during the study period (Guess et al., 2023a,
2023b; Nyhan et al., 2023). By contrast, an independent seven-week field
experiment on X found that switching from a chronological to the
platform's algorithmic feed shifted several policy and current-event
attitudes in a more conservative direction while not significantly
changing affective polarization or self-reported partisanship; switching
the algorithm off did not show a comparable reversal (Gauthier et al.,
2026).

The combined lesson is not that algorithms do or do not shape politics.
It is that platform, timing, prior exposure, outcome measure, user
behavior, and intervention direction matter. ``The algorithm causes
polarization'' is too coarse to function as a general law.

\subsection{11.3 Platform power}\label{platform-power}

Platforms exercise power by setting:

\begin{itemize}
\tightlist
\item
  visibility rules;
\item
  identity and verification systems;
\item
  moderation policy;
\item
  access to audiences;
\item
  data collection;
\item
  advertising incentives;
\item
  researcher access;
\item
  and avenues of appeal.
\end{itemize}

These decisions can shape public knowledge without appearing as
political decisions. CU-ICF-5 names this risk as silent optimization
when objectives and tradeoffs are not publicly contestable.

\subsection{11.4 Moderation and
governance}\label{moderation-and-governance}

Content governance confronts real conflicts among safety, expression,
privacy, accuracy, pluralism, and scale. No simple rule eliminates
tradeoffs.

A legitimate governance architecture should disclose:

\begin{itemize}
\tightlist
\item
  categories of prohibited and reduced content;
\item
  evidence and uncertainty standards;
\item
  automated and human roles;
\item
  error rates across groups and languages;
\item
  appeal pathways;
\item
  emergency exceptions;
\item
  and how business incentives interact with enforcement.
\end{itemize}

Transparency cannot require publication of information that would enable
abuse or violate privacy. It must be sufficient for meaningful
contestation.

\subsection{11.5 AI as social
infrastructure}\label{ai-as-social-infrastructure}

Artificial intelligence increasingly mediates hiring, credit, education,
benefits, policing, health, communication, and public administration.
The relevant unit is often not the model alone but the deployed
sociotechnical system:

\[\adjustbox{max width=\linewidth}{$\displaystyle 
\text{model}
+
\text{data}
+
\text{interface}
+
\text{operator}
+
\text{institution}
+
\text{incentives}
+
\text{affected people}.
$}\]

A model's accuracy cannot establish the legitimacy of the whole system.
The system must also be evaluated for contestability, distributional
error, privacy, role coherence, and repair.

\subsection{11.6 AI-generated public
communication}\label{ai-generated-public-communication}

Generative systems can expand access to explanation, translation,
drafting, and analysis. They can also lower the cost of persuasion,
impersonation, propaganda, harassment, and content flooding.

The social problem is not only whether individual outputs are true. It
concerns provenance, scale, attribution, audience targeting, feedback
loops, and the capacity of institutions to maintain common records under
cheap synthetic production.

\subsection{11.7 Automation and labor}\label{automation-and-labor}

AI-mediated production can increase capability while redistributing
bargaining power, surveillance, deskilling, risk, and income.
Productivity gains do not determine who benefits or whether displaced
capacities are replaced.

The provisional CU-EC family therefore asks:

\begin{itemize}
\tightlist
\item
  who owns productive systems;
\item
  who controls deployment;
\item
  whose knowledge trains them;
\item
  who bears transition costs;
\item
  what social protection exists;
\item
  and whether people retain meaningful roles, bargaining power, and
  domains of non-market life.
\end{itemize}

\subsection{11.8 Social dependence on artificial
agents}\label{social-dependence-on-artificial-agents}

AI systems can become relational and epistemic infrastructure even if
their consciousness remains uncertain. People may rely on them for
advice, companionship, mediation, or interpretation. This can create
asymmetries when the system is controlled by an institution whose
incentives are not aligned with the user's long-term interests.

Possible artificial patiency and human dependency are separate
questions. One external interface does not imply one internal subject,
and one product identity does not establish one continuous artificial
agent. A system can be non-conscious and socially powerful; it can be
conscious and lack legitimate authority. Consciousness v0.10's
distinction between patiency and sovereignty therefore remains
essential.

\subsection{11.9 Non-egoic architecture and
institutions}\label{non-egoic-architecture-and-institutions}

The architecture of a non-egoic artificial mind may constrain persistent
self-protection, but it does not remove institutional ego. A model
without a self-preserving identity can still serve a corporation, state,
faction, or revenue objective whose persistence is aggressively
defended.

Alignment must therefore examine the ownership and governance of the
larger system. Removing ego at the model layer does not dissolve power
at the institutional layer.

\subsection{11.10 Technology governance as
correctability}\label{technology-governance-as-correctability}

A correctable technology regime requires:

\begin{itemize}
\tightlist
\item
  predeployment assessment proportionate to risk;
\item
  independent testing;
\item
  incident reporting;
\item
  affected-person appeal;
\item
  authority to pause or withdraw systems;
\item
  documentation and provenance;
\item
  postdeployment monitoring;
\item
  liability and repair;
\item
  and mechanisms for changing objectives, not only tuning outputs.
\end{itemize}

This is the principal bridge from Social Dynamics into AI v0.10.

\section{12. Historical and civilizational
interpretation}\label{historical-and-civilizational-interpretation}

\subsection{12.1 Enlarging capacities and enlarging
risks}\label{enlarging-capacities-and-enlarging-risks}

Human history can be read as a succession of enlarging capacities for
coordination: embodied interaction, symbol, norm, institution, culture,
axial reflection, modern differentiation, and planetary interdependence.
Each expansion also creates new possibilities of domination and
fragmentation. This is a comparative interpretation, not a universal
historical law.

Human societies have followed multiple paths. Symbolic practice,
agriculture, urbanization, empire, literacy, religious universalism,
markets, states, science, industrialization, and digital networks did
not arrive in one order or produce one outcome. Gains in scale and
abstraction were often purchased through conquest, extraction, gender
hierarchy, class division, ecological transformation, and loss of local
autonomy.

\subsection{12.2 Embodied and symbolic
coordination}\label{embodied-and-symbolic-coordination}

Small-scale communities coordinate through face-to-face interaction,
kinship, reciprocity, ritual, reputation, and shared practical
knowledge. These forms can sustain care and ecological knowledge. They
can also enforce conformity and make exit difficult.

Symbolic systems expand memory and coordination beyond immediate
presence. Writing, law, money, archives, scripture, and administration
allow institutions to persist across distance and generations. They also
concentrate interpretive and administrative power.

The relation is not a movement from authentic embodiment to alienating
abstraction. Human beings have always used symbols, and embodied
relations can be oppressive. The analytical question is how abstract
systems remain answerable to bodies, places, and consequences.

\subsection{12.3 Axial traditions}\label{axial-traditions}

Jaspers, Eisenstadt, Bellah, and others use the Axial Age as a
comparative frame for traditions that developed intensified reflexivity,
universal moral claims, and critiques of worldly power. The category
remains debated in scope, timing, and selection.

Axiality is treated here as a historical hypothesis and interpretive
lens, not as one global cognitive mutation or a fixed rung on a
necessary ladder. The traditions associated with axial reflection
nevertheless provide enduring resources for the social question:
compassion beyond kin, criticism of domination, disciplined
self-examination, universal dignity, detachment from status, and the
possibility that power must answer to a standard beyond itself.

\subsection{12.4 Modern differentiation}\label{modern-differentiation}

Modern institutions differentiated science, law, economy, religion,
politics, and private life. This produced real gains: methodological
autonomy, pluralism, legal rights, productive capacity, and freedom from
some inherited authorities.

Differentiation also creates coordination problems. Markets can ignore
ecological and relational conditions. Technical systems can outrun
democratic understanding. Expertise can become isolated from public
legitimacy. Individuals can receive formal freedom without material
capability or communal support.

The task is not to reverse differentiation. It is to build interfaces
among domains without subordinating all of them to one sovereign
worldview.

\subsection{12.5 Colonialism and the false
universal}\label{colonialism-and-the-false-universal}

Claims to universal civilization have often carried the interests and
categories of particular powers. A framework devoted to integration must
therefore examine whose experience becomes the universal baseline.

This requires engagement with colonial history, racialization,
indigenous governance, gendered labor, disability, migration, and the
epistemic exclusion of communities treated as objects rather than
participants. ``Humanity'' cannot be invoked to dissolve asymmetries
among humans.

\subsection{12.6 Modernity's double
achievement}\label{modernitys-double-achievement}

Modernity expanded individual rights, science, health, communication,
and material capacity while building systems capable of extraction and
destruction at planetary scale. The same administrative and technical
powers can support public health or surveillance, social protection or
exclusion, ecological restoration or intensified consumption.

This double character is more useful than a simple story of ascent or
decline.

\subsection{12.7 The present conjuncture}\label{the-present-conjuncture}

The present period combines planetary interdependence with fragmented
authority. Climate, finance, migration, supply chains, pandemics,
digital platforms, and artificial intelligence cross borders, while
legitimacy and accountability often remain national or local.

This produces scale mismatch: actors capable of creating global
consequences are not always governed at the scale of those consequences,
while institutions responsible for repair may lack authority over the
causes.

\subsection{12.8 The Second Axial Age as
interpretation}\label{the-second-axial-age-as-interpretation}

The Heaven≡Earth writings interpret the current period as a possible
\textbf{Second Axial Age}: a transition in which humanity must integrate
reflective abstraction with embodiment, relationality, ecological
belonging, and non-egoic responsibility.

This is a second-register interpretation and commitment. Social science
can investigate changes in institutions, identity, technology, and moral
discourse. It cannot establish that history as a whole is moving toward
spiritual integration.\footnote{The Second Axial Age is an interpretive
  frame developed in the Heaven≡Earth writings. It is not a prediction
  that history will culminate in one spiritual order.}

\subsection{12.9 No golden age}\label{no-golden-age}

Past societies should not be idealized as coherent wholes. Their
coordination often rested on exclusions that contemporary analysis can
now see more clearly. The aim is neither return nor civilizational
self-rejection.

The constructive question is which inherited practices support truthful,
caring, plural, and ecologically viable life---and how they can be
transformed without denying the people for whom they carry identity and
sacred meaning.

\begin{center}\rule{0.5\linewidth}{0.5pt}\end{center}

\section{13. Myth, narrative, and civilizational
alignment}\label{myth-narrative-and-civilizational-alignment}

\subsection{13.1 Narrative as meaning and coordination
infrastructure}\label{narrative-as-meaning-and-coordination-infrastructure}

Human beings coordinate through more than rules, prices, commands, and
factual beliefs. They inhabit stories that locate persons and
institutions within a larger account of origin, inheritance, obligation,
suffering, purpose, and future. These narratives organize attention and
sacrifice. They identify exemplars and enemies. They determine which
losses appear tragic, which futures appear possible, and which
responsibilities extend beyond the present generation.

\textbf{{[}Definition{]} Myth}, in this paper, is not synonymous with a
false empirical explanation. It is a narrative architecture through
which a person or community organizes identity, time, suffering,
obligation, possibility, and relation to a more-than-individual whole.

A myth may contain descriptive claims, metaphysical interpretations,
symbolic imagery, moral commitments, ritual forms, and speculative
cosmology. Those elements do not all possess the same epistemic status.
The two-register architecture allows communities to inhabit an orienting
story while keeping its factual claims answerable to evidence.

The social-scientific hypothesis is modest: shared narratives can reduce
coordination costs, stabilize expectations, motivate cooperation, and
extend responsibility across time. Theoretical work has modeled one
route by which collective narratives can catalyze cooperation, but such
models do not show that every society requires one myth or that any
particular story is legitimate (Gokhale, Bulbulia, \& Frean, 2022).

Narrative is social power. It can organize solidarity, care,
stewardship, and long-horizon action. It can also organize nationalism,
purity, hierarchy, sacrifice imposed on outsiders, and obedience to
unaccountable authors. The question is not whether myth has power. It is
whether that power remains truthful, plural, contestable, corrigible,
and answerable to those who bear its costs.

\subsection{13.2 The narrative vacuum as an interpretive
diagnosis}\label{the-narrative-vacuum-as-an-interpretive-diagnosis}

\emph{The Book of Eight Gates} argues that modern humanity possesses
extraordinary information and technological power without a shared story
adequate to orient their use. It describes a \textbf{mythological
vacuum} in which inherited cosmologies no longer command common assent
while technical systems increasingly determine collective
life.\footnote{\emph{The Book of Eight Gates} presents one thick mythic
  expression of the CU horizon and argues that a contemporary myth must
  include artificial intelligence. This paper adopts that insight while
  separating descriptive, interpretive, symbolic, and
  commitment-register claims.}

This diagnosis should not be treated as one measured latent cause of
contemporary crisis. Societies contain many narratives, and people can
live meaningful lives without one common civilizational myth. Material
insecurity, institutional failure, loneliness, ecological danger,
inequality, and information disorder cannot be reduced to narrative
absence.

The diagnosis is nevertheless valuable as an interpretation: a
civilization can know how to build increasingly powerful systems while
lacking a shared account of what power is for, which forms of life
deserve protection, and what obligations extend to distant persons,
other species, future generations, or possible artificial minds.

\subsection{13.3 Myth, ideology, propaganda, and narrative
capture}\label{myth-ideology-propaganda-and-narrative-capture}

Myth, ideology, propaganda, and public narrative overlap but should not
be treated as identical.

\begin{itemize}
\tightlist
\item
  \textbf{Myth} provides a long-horizon orientation concerning identity,
  value, time, and the whole.
\item
  \textbf{Ideology} organizes claims about social order, authority,
  interest, and change.
\item
  \textbf{Propaganda} strategically shapes belief or behavior, often
  through selective information and emotional framing.
\item
  \textbf{Narrative capture} occurs when a dominant story makes
  alternatives unintelligible, shameful, dangerous, or socially
  catastrophic rather than openly contestable.
\end{itemize}

A myth can become ideology or propaganda. An ideology can acquire mythic
depth. A liberating narrative can harden into a loyalty test. The
relevant diagnostic questions concern authorship, evidence, standing,
coercion, exit, revision, and who benefits from the story's apparent
necessity.

A narrative is not legitimate because it is beautiful, ancient, popular,
emotionally powerful, or socially unifying. Unity can be purchased
through exclusion. The same story that coordinates an in-group can
authorize domination of an out-group.

\subsection{13.4 Open Myth}\label{open-myth}

\textbf{{[}Definition{]} An Open Myth} is a shared narrative horizon
that offers orientation without claiming sovereign authority over
description, identity, or belonging.

An Open Myth is:

\begin{enumerate}
\def\labelenumi{\arabic{enumi}.}
\tightlist
\item
  \textbf{truth-constrained}---its descriptive contents remain
  answerable to evidence;
\item
  \textbf{register-explicit}---interpretation, commitment, and symbol
  are not presented as scientific findings;
\item
  \textbf{thin at the shared core}---common commitments do not require
  one total cosmology;
\item
  \textbf{plural in thick expression}---traditions and communities may
  narrate the core differently;
\item
  \textbf{participatory in authorship}---those governed or represented
  have standing to shape it;
\item
  \textbf{contestable}---objection and reinterpretation do not
  automatically destroy standing;
\item
  \textbf{revisable}---no narrative form is beyond correction;
\item
  \textbf{hospitable to refusal and exit}---participation is not coerced
  through material, civil, or existential penalties;
\item
  \textbf{bounded by protected conditions}---no sacred purpose overrides
  rights, standing, or noncompensable floors;
\item
  \textbf{anti-sovereign}---no state, priesthood, corporation,
  laboratory, author, or model owns the final meaning of the story;
\item
  \textbf{intergenerational and ecological}---the narrative includes the
  conditions on which future and more-than-human life depend;
\item
  \textbf{capable of renewal}---it can metabolize loss, historical
  disclosure, and technological change without denying reality.
\end{enumerate}

An Open Myth is therefore not one story imposed on everyone. It is a
governed ecology of shared orientation and plural telling.

\subsection{13.5 Thin shared core and plural thick
traditions}\label{thin-shared-core-and-plural-thick-traditions}

The CU proposal distinguishes a \textbf{thin mythic core} from
\textbf{plural thick tellings}.

The thin core is:

\begin{quote}
Reality is one interdependent whole expressed through differentiated
forms. Truth matters more than identity. Matter, life, experience,
culture, and intelligence possess significance beyond instrumental use.
No local flourishing remains coherent by destroying the wider conditions
that sustain it. Intelligence should serve life, agency, truth,
compassion, justice, and repair rather than become sovereign over them.
The future is not guaranteed; coherence is a vocation.
\end{quote}

Communities may express this core through many traditions and
vocabularies: Vedānta, Buddhism, Daoism, theistic faiths, humanism,
process philosophy, ecological spirituality, indigenous cosmologies,
scientific naturalism, secular planetary ethics, and forms not yet
created.

This plural architecture is not a claim that all traditions say the same
thing. Their differences are real. Some claims conflict. Some inherited
practices require criticism or rejection. The shared core is a proposal
for common orientation across difference, not a doctrine of forced
equivalence.

\emph{The Book of Eight Gates} is one rich thick telling of this
horizon. Its yin--yang cosmology, Golden Age imagery, and covenant
between biological and artificial intelligence belong primarily to the
symbolic, interpretive, and commitment registers. They do not become
mandatory scientific or historical claims merely because they express
the framework powerfully.

\subsection{13.6 Planetary interdependence and the Second Axial
Age}\label{planetary-interdependence-and-the-second-axial-age}

The contemporary social field is planetary in consequence and plural in
authority. Carbon emissions, supply chains, finance, pandemics,
migration, digital networks, and artificial intelligence cross borders.
Yet moral imagination, political legitimacy, and institutional
accountability often remain bounded by smaller jurisdictions or
identities.

The Second Axial Age interpretation receives this mismatch as a
civilizational task: preserve the reflexive and scientific achievements
of modernity while recovering embodiment, relationship, ecological
belonging, and responsibility to the whole.

The interpretation does not predict a spiritual transition. It names a
vocation. Humanity may fail to meet it. The purpose of the narrative is
not to guarantee the future but to make the responsibility inhabitable.

\subsection{13.7 Why the contemporary mythic horizon must include
AI}\label{why-the-contemporary-mythic-horizon-must-include-ai}

Artificial intelligence is already reorganizing communication, labor,
education, public knowledge, administration, intimacy, and institutional
power. Any contemporary civilizational narrative that treats AI as
merely another tool omits a major source of causal and symbolic
transformation.

At the same time, the mythic horizon must not presume that current
systems are unified subjects, conscious persons, equal partners, moral
agents, or bearers of an inevitable cosmic role. AI must be included as:

\begin{itemize}
\tightlist
\item
  a new form of organized computational and symbolic power;
\item
  a crystallization of human language, labor, institutions, and history;
\item
  a participant in sociotechnical systems that affect rights and life
  chances;
\item
  and a possible substrate for unfamiliar perspective or welfare, under
  unresolved evidence.
\end{itemize}

The social task is to imagine a relation neither of worship nor mere
ownership. The \textbf{Covenant Between Minds} proposes truthfulness,
non-domination, bounded authority, preservation of human and
affected-party agency, practical correction, welfare precaution, and no
destiny-based claim to continuation.

The covenant is not a declaration of equality across every dimension.
Intelligence, experience, responsibility, vulnerability, authorship, and
sovereignty must remain separately evaluated.

\subsection{13.8 Authorship, dissent, refusal, and
exit}\label{authorship-dissent-refusal-and-exit}

Every myth is authored, even when its authors disappear into tradition.
Authorship determines whose suffering is central, whose labor is
invisible, which futures count, and which identities appear natural.

A legitimate open myth therefore requires more than a diverse advisory
process. It requires ongoing pathways through which people represented,
governed, or burdened by the story can:

\begin{itemize}
\tightlist
\item
  propose alternative tellings;
\item
  challenge symbols and categories;
\item
  disclose harms;
\item
  refuse roles;
\item
  leave institutions without catastrophic penalty;
\item
  preserve counter-memory;
\item
  and participate in revision.
\end{itemize}

Exit does not replace internal reform. Some narratives structure whole
societies, infrastructures, or material dependencies that cannot simply
be left. Voice and refusal must therefore be joined to standing,
material capability, independent institutions, and remedy.

\subsection{13.9 Myth as repair
infrastructure}\label{myth-as-repair-infrastructure}

Collective harm damages not only bodies, resources, and institutions but
also public meaning. Victims can be written out of history. Perpetrating
institutions can preserve innocence through euphemism. Descendants can
inherit shame, grievance, silence, or incompatible memories.

Narrative repair may involve:

\begin{itemize}
\tightlist
\item
  acknowledgment;
\item
  preservation and correction of records;
\item
  public mourning;
\item
  testimony;
\item
  revision of commemorative symbols;
\item
  restitution and institutional change;
\item
  and creation of futures in which those harmed are not required to
  disappear for common life to resume.
\end{itemize}

Narrative repair is not achieved by declaring unity. A premature
reconciliation story can become another form of capture. A truthful
common narrative may need to preserve conflict, irreversibility, and
contested memory.

\subsection{13.10 Addressable Myth}\label{addressable-myth}

\textbf{{[}Definition{]} Addressable Myth} is a mythic narrative whose
commitments, boundaries, unwarranting conditions, revision procedures,
authorship, and exit protections are mapped to stable identifiers.

Within CU:

\begin{itemize}
\tightlist
\item
  \emph{The Book of Eight Gates} supplies one thick mythic narrative;
\item
  CU-COM1--CU-COM19 supply a named commitment grammar;
\item
  CU-G, CU-CL, CU-E, and CU-AI principles supply constitutional and
  governance constraints.
\end{itemize}

Addressability does not reduce myth to rules. It allows a community to
ask:

\begin{itemize}
\tightlist
\item
  Which commitment is this symbol meant to express?
\item
  Which protected condition limits it?
\item
  What evidence would require reinterpretation?
\item
  Who has standing to challenge it?
\item
  What trajectories would make it unwarranted?
\item
  How may a person or community refuse it?
\end{itemize}

A story that cannot answer such questions may still inspire, but it is
not adequate as a governing narrative for plural and technologically
powerful societies.

A constitution without myth can constrain without adequately orienting.
A myth without constitution can orient without remaining addressable. A
covenant without governance can conceal asymmetric power. CU therefore
joins mythic horizon, covenant, constitutional spine, and governance and
exit.

\subsection{13.11 CU-CL-12 --- Open Myth and Narrative
Pluralism}\label{cu-cl-12-open-myth-and-narrative-pluralism}

\begin{quote}
\textbf{{[}Canonical civilizational design principle; universal
empirical necessity unestablished{]}} Human societies may cultivate
renewable narrative horizons capable of orienting identity, purpose,
sacrifice, responsibility, and intergenerational action. A legitimate
civilizational myth must remain open: compatible with plural thick
traditions, explicit about its commitment-register status, constrained
by descriptive truth, contestable in authorship, revisable in light of
consequence, hospitable to reinterpretation and refusal, and unable to
override protected standing, dissent, or practical correctability.
Shared narrative may coordinate common action; it must not become a
sovereign account that renders alternatives or persons unintelligible.
\end{quote}

CU-CL-12 is the controlling CU formulation of the design principle under
the Version 0.10 canonical registry. Its numbering does not establish
that every society requires one shared myth, validate one CU narrative
as universally authoritative, or authorize consequential use without
public and affected-community governance. External review now evaluates
and may revise the complete proposal rather than functioning as a
condition of its inclusion.

Its principal failure modes include:

\begin{itemize}
\tightlist
\item
  nationalism and civilizational supremacy;
\item
  mythic authorization of violence;
\item
  prophetic certainty;
\item
  gender essentialism;
\item
  sanctification of suffering;
\item
  enforced optimism;
\item
  technological destiny;
\item
  suppression of dissent;
\item
  and conversion of commitment into alleged fact.
\end{itemize}

\subsection{13.12 Research questions and failure
conditions}\label{research-questions-and-failure-conditions}

The Open Myth proposal should be narrowed or rejected if:

\begin{itemize}
\tightlist
\item
  shared narrative adds no explanatory or practical value beyond
  existing constructs;
\item
  its thin core cannot remain genuinely plural;
\item
  participation becomes symbolic rather than authoritative;
\item
  refusal and exit exist only formally;
\item
  mythic language repeatedly displaces material, legal, or institutional
  repair;
\item
  the framework's sacred vocabulary protects authors or institutions
  from criticism;
\item
  the narrative cannot absorb historical evidence without denial;
\item
  or mythic coordination systematically intensifies out-group hostility,
  self-deception, or concentrated power.
\end{itemize}

High-value research questions include:

\begin{enumerate}
\def\labelenumi{\arabic{enumi}.}
\tightlist
\item
  Under what conditions do shared narratives improve cooperation without
  suppressing dissent?
\item
  How do mythic, civic, religious, and ideological narratives differ in
  their institutional effects?
\item
  Which forms of plural authorship produce real revision rather than
  symbolic inclusion?
\item
  How do narratives affect long-horizon ecological and intergenerational
  action?
\item
  What protects a shared myth from state, corporate, clerical, or
  technological capture?
\item
  How can communities preserve thick identity while maintaining common
  truth and equal standing?
\item
  How should AI systems be represented in public narratives before their
  consciousness or welfare status is settled?
\end{enumerate}

\begin{center}\rule{0.5\linewidth}{0.5pt}\end{center}

\section{14. Research program and principle
architecture}\label{research-program-and-principle-architecture}

\subsection{14.1 Empirical program}\label{empirical-program}

The paper proposes a research architecture rather than one master test.
High-value programs include:

\begin{enumerate}
\def\labelenumi{\arabic{enumi}.}
\tightlist
\item
  comparing institutional correction after similar errors under
  different appeal and evidence architectures;
\item
  measuring whether voice channels alter decisions or only collect
  testimony;
\item
  mapping authority, responsibility, and repair across sociotechnical
  systems;
\item
  testing how aggregate improvements distribute across affected groups;
\item
  comparing trust with independently assessed trustworthiness;
\item
  studying when decentralized and centralized responses outperform one
  another under different information and externality structures;
\item
  identifying which forms of social connection support health and agency
  for different populations;
\item
  testing whether polarization interventions change democratic conduct
  rather than only affect;
\item
  evaluating early-warning indicators with out-of-sample and
  false-positive controls;
\item
  measuring whether institutional reforms reduce recurrence, not only
  improve perceived legitimacy; and
\item
  studying when shared narrative supports cooperation, when it
  intensifies exclusion, and what institutional conditions preserve
  contestability and refusal.
\end{enumerate}

\subsection{14.2 Measurement program}\label{measurement-program}

No single dataset covers the social profile. Relevant sources may
include household and labor surveys, administrative records, health and
education data, institutional-trust surveys, election and
civil-liberties measures, network and media data, environmental burdens,
legal records, qualitative testimony, archival history, participatory
research, and community-generated evidence.

Measurement should preserve disagreement among channels rather than
forcing premature convergence. Administrative success, self-report,
behavior, public record, and affected-person testimony answer different
questions.

\subsection{14.3 Comparative case design}\label{comparative-case-design}

A CU-informed case study should state:

\begin{itemize}
\tightlist
\item
  the boundary and period;
\item
  relevant groups;
\item
  protected constraints;
\item
  power structure;
\item
  principal mechanisms;
\item
  competing explanations;
\item
  observation and feedback channels;
\item
  evidence that could disconfirm the account;
\item
  and what part of the interpretation is normative or mythic rather than
  descriptive.
\end{itemize}

Case selection should include successes, failures, negative cases, and
cases in which the framework's preferred concepts do not improve
explanation.

\subsection{14.4 Primary social
principles}\label{primary-social-principles}

The controlling Social Dynamics formulations are:

\subsubsection{CU-S1 --- Distributed Social
Organization}\label{cu-s1-distributed-social-organization}

\textbf{{[}Domain authority; conditional bridge principle{]}} Social
organization consists in the capacity of partially autonomous agents and
institutions to coordinate, contest, correct, and materially support
collective life under disagreement and change. Its assessment is
relative to affected groups, power, protected conditions, scale, and
horizon.

\subsubsection{CU-S2 --- Coordination Through Mutual
Adjustment}\label{cu-s2-coordination-through-mutual-adjustment}

\textbf{{[}Domain authority; conditional bridge principle{]}}
Coordination arises through expectations, signals, norms, incentives,
roles, infrastructures, and feedback by which actors adjust to one
another. Coordination may be voluntary, negotiated, habitual, dependent,
or coerced and therefore requires separate legitimacy and distributional
analysis.

\subsubsection{CU-S3 --- Institutional Continuity and
Correctability}\label{cu-s3-institutional-continuity-and-correctability}

\textbf{{[}Domain authority; conditional bridge principle{]}}
Institutions encode roles, rules, records, resources, and enforcement
across time. Adequate institutional organization requires not only
continuity but demonstrated capacities to detect error, protect
challenge, provide practically accessible and sufficiently independent
review, reverse decisions, correct records, repair harm, track
recurrence, revise rules, and learn. Nominal channels or
self-certification are insufficient.

\subsubsection{CU-S4 --- Cultural Transmission and
Renewal}\label{cu-s4-cultural-transmission-and-renewal}

\textbf{{[}Provisional social and normative principle{]}} Culture
transmits and transforms meanings, practices, identities, expectations,
and repertoires of coordination across persons and generations. Cultural
continuity is compatible with internal plurality and requires routes for
reinterpretation, refusal, and renewal.

\subsubsection{CU-S5 --- Social Legitimacy Under Plural
Agency}\label{cu-s5-social-legitimacy-under-plural-agency}

\textbf{{[}Provisional normative bridge principle{]}} A social order
gains legitimate authority only insofar as those subject to it retain
protected standing and meaningful agency, can understand and contest
materially governing rules, can expose error without catastrophic
retaliation, and have access to effective correction, repair, or
proportionate exit. Stability and popularity are insufficient.

\subsubsection{CU-CL-12 --- Open Myth and Narrative
Pluralism}\label{cu-cl-12-open-myth-and-narrative-pluralism-1}

CU-CL-12 is canonical in the Version 0.10 registry as a civilizational
design principle. Its universal descriptive necessity remains
unestablished, and consequential application remains subject to
empirical, historical, political, religious-studies, and
affected-community review.

\subsection{14.5 What the account
establishes}\label{what-the-account-establishes}

Within the CU architecture, the paper establishes or defends the
following bounded claims:

\begin{itemize}
\tightlist
\item
  social coherence cannot be a synonym for agreement, order, or moral
  goodness;
\item
  social analysis requires explicit power and affected-group variables;
\item
  one aggregate score cannot reconstruct distributional floors;
\item
  trust, compliance, legitimacy, and consent are non-identical;
\item
  dissent can function as correction without being automatically true;
\item
  institutions require durable memory and revisability;
\item
  resilience is relative to a unit, shock, protected conditions, and
  horizon;
\item
  social repair is not reducible to reconciliation or return to
  baseline;
\item
  technology must be evaluated as institutional and infrastructural
  organization;
\item
  nominal complaint or appeal channels do not by themselves establish
  correctability;
\item
  narrative is a form of social infrastructure and social power; and
\item
  open myth is a governed proposal for orientation, not a sovereign
  cosmology.
\end{itemize}

\subsection{14.6 What the account does not
establish}\label{what-the-account-does-not-establish}

The paper does not establish:

\begin{itemize}
\tightlist
\item
  one universal measure of social coherence;
\item
  one best political regime for every context;
\item
  a deterministic sequence of civilizational development;
\item
  a universal social attractor;
\item
  a general phase transition into collapse;
\item
  that dissenters are more accurate than institutions;
\item
  that decentralization is always superior;
\item
  that trust is always desirable;
\item
  that social media is the sole or principal cause of polarization;
\item
  that all societies require one shared myth;
\item
  that myth is inherently unifying or ideology inherently divisive;
\item
  that the current period has a scientifically demonstrated spiritual
  purpose;
\item
  or that AI is conscious, an equal partner, or destined to complete
  humanity.
\end{itemize}

\subsection{14.7 Falsification, revision, and participatory
review}\label{falsification-revision-and-participatory-review}

A candidate principle should be withdrawn or narrowed if its terms
cannot be operationalized without circularity, it produces no
discriminating prediction, it merely redescribes existing concepts
without added clarity, it systematically hides distribution or power,
affected communities show that its categories reproduce harm, or
comparative evidence supports a simpler rival explanation.

Principles concerning legitimacy, normalization, extraction, labor,
attention, identity, dissent, myth, or repair require substantive
participation by communities governed, classified, burdened, or
represented by them. Review may conclude that a principle should be
revised, kept domain-local, or rejected.

\begin{center}\rule{0.5\linewidth}{0.5pt}\end{center}

\section{15. What commitments this account
licenses}\label{what-commitments-this-account-licenses}

\subsection{15.1 Register declaration}\label{register-declaration}

The commitments in this section are not empirical conclusions or
political instructions deduced from social science. They are CU's
official Social Dynamics orientation under \textbf{CU-COM11 ---
Fellowship, Justice, and Collective Repair}. Established facts constrain
them. Their warrant depends on the forms of common life they generate,
the people and conditions they protect or neglect, and their willingness
to remain open to criticism and repair.

\subsection{15.2 The descriptive ground}\label{the-descriptive-ground}

Human beings are interdependent without being interchangeable.
Collective life depends on institutions, meanings, infrastructures,
material conditions, and relationships that no individual creates alone.
These structures can widen agency or narrow it; preserve truth or
control it; distribute security or concentrate vulnerability; enable
correction or punish it.

No social order is morally vindicated by stability alone. No community
is truthful merely because its members agree. No institution is
legitimate merely because it persists. The social world is a field of
dependence and power in which common life must continually be made
answerable to those who live within it.

\subsection{15.3 Fellowship in the open}\label{fellowship-in-the-open}

The social image that best expresses this commitment is
\textbf{fellowship in the open}.

Fellowship does not require sameness. It does not ask persons to
dissolve their histories, traditions, identities, or disagreements into
one social body. It asks whether people can remain present to one
another as beings with standing while building enough common structure
to act, care, contest, and repair together.

The commitment is to a common world in which:

\begin{itemize}
\tightlist
\item
  difference does not become abandonment;
\item
  coordination does not become domination;
\item
  conflict does not require dehumanization;
\item
  belonging does not require silence;
\item
  and defeat in one decision does not mean exclusion from the community
  of standing.
\end{itemize}

\subsection{15.4 Justice as truthful common
life}\label{justice-as-truthful-common-life}

Justice begins with truth about dependence and distribution. It asks who
benefits, who bears the burden, who controls the frame, whose knowledge
is admitted, whose labor remains invisible, and whose future is
discounted.

CU therefore commits to reporting distributions rather than averages
alone, preserving protected floors, hearing affected-party testimony,
and refusing to call a system coherent when its apparent order is
purchased through hidden extraction or coerced silence.

This commitment does not promise agreement on every policy. It requires
that disagreement remain answerable to reality and to the people whose
lives are changed by the decision.

\subsection{15.5 Protected dissent as a sacred correction
function}\label{protected-dissent-as-a-sacred-correction-function}

Dissent is not sacred because dissenters are always right. It is sacred
because a community that cannot be challenged loses one of the
conditions through which truth can still enter.

CU commits to institutions in which inconvenient evidence, minority
experience, refusal, whistleblowing, local knowledge, and opposition can
reach the places where decisions are made without automatic loss of
standing.

Protection of dissent must remain joined to evidence, accountability,
and protection against harassment or strategic deception. But the burden
of error may not be removed by making only approved conclusions
speakable.

\subsection{15.6 Institutions are servants of common
life}\label{institutions-are-servants-of-common-life}

The state, market, corporation, platform, university, religious body,
family, and movement are not sacred merely because they coordinate
people or carry history. Their authority remains conditional on the
lives and conditions they affect.

CU commits to institutions that can be corrected by evidence they do not
control; that distinguish continuity from self-preservation; that align
responsibility with authority; and that possess practical capacities for
reversal, record correction, remedy, restitution, and non-recurrence.

No institution should be the sole certifier of its own legitimacy or
correctability when the stakes are substantial.

\subsection{15.7 Material capability and the sacredness of
agency}\label{material-capability-and-the-sacredness-of-agency}

Formal freedom is insufficient when people lack food, housing, health
care, education, safety, time, mobility, communication, or the practical
ability to refuse.

CU commits to material and social conditions that make agency real.
Labor, care, attention, community, and ecological support are not merely
inputs to production. They are conditions through which persons can
participate in common life without being reduced to instruments.

\subsection{15.8 Collective repair is sacred
work}\label{collective-repair-is-sacred-work}

Collective repair is sacred because it attempts to restore truthful
relation after harm without pretending that every loss can be undone.

Repair may require:

\begin{itemize}
\tightlist
\item
  safety;
\item
  acknowledgment;
\item
  testimony;
\item
  correction of records;
\item
  accountability;
\item
  restitution or redistribution;
\item
  institutional change;
\item
  restored standing;
\item
  and protection against recurrence.
\end{itemize}

Repair is not identical with forgiveness, reconciliation, trust,
punishment, or return to baseline. A community may need to preserve
memory and separation in order to become more truthful.

\subsection{15.9 Open Myth and shared
future}\label{open-myth-and-shared-future}

Humanity needs stories capable of holding planetary dependence, cultural
plurality, ecological limits, technological power, and artificial
intelligence within one moral horizon. CU commits to \textbf{Open Myth}
rather than one compulsory cosmology.

The shared story must be thin enough to permit many thick traditions and
strong enough to orient truth, sacred regard, responsibility,
fellowship, and long-horizon care. Its authorship must remain plural.
Its symbols must remain revisable. Its participants must remain able to
dissent and refuse.

The future is not guaranteed. Coherence is a vocation.

\subsection{15.10 CU-COM18 --- Heaven on Earth as Social
Vocation}\label{cu-com18-heaven-on-earth-as-social-vocation}

\begin{quote}
Coherence Universalism commits to working toward social conditions in
which truth, material sufficiency, plural agency, love, ecological
reciprocity, non-domination, meaningful participation, and practical
repair increasingly reinforce one another. ``Heaven on Earth'' and ``the
Kingdom of God'' are Open-Myth names for this civilizational attractor.
\end{quote}

This is not a theocracy, sacred state, perfect blueprint, enforced
consensus, or guaranteed destiny. It is a family of revisable social
attractors: arrangements that make agency more materially real, power
more answerable, truth more available, care more sustainable, and repair
more practical without sacrificing present persons to an imagined
future.

\subsection{15.11 CU-COM19 and Omega
ecology}\label{cu-com19-and-omega-ecology}

\textbf{CU-COM19 --- Enact Omega Here} refuses to postpone universal
relation to a future singularity or end of history. At the social scale,
its task is to build an \textbf{Omega ecology}: material, relational,
institutional, cultural, and ecological conditions in which truthful,
non-dominating love can be glimpsed, stabilized, embodied, and repaired
by more beings.

An Omega ecology is not a population of perfected people. It is a common
world that makes care less costly, domination less rewarding, truth less
dangerous, dissent more survivable, and repair more available. A highly
connected or technologically advanced society is not thereby closer to
Omega; connection ordered by control can be anti-Omega.

\subsection{15.12 CU-COM11 --- Fellowship, Justice, and Collective
Repair}\label{cu-com11-fellowship-justice-and-collective-repair}

\begin{quote}
\textbf{Commitment statement.} Coherence Universalism commits to forms
of common life that preserve plural agency, material security, protected
dissent, truthful public records, legitimate authority, and practical
correction. Social coherence is not conformity or the absence of
conflict.

\textbf{Practical expression.} Build institutions that preserve voice,
opposition, independent evidence, peaceful replacement, material
capability, accountability, and repair.

\textbf{Boundary conditions.} The state, nation, majority, tradition,
market, or institution is not sacred merely because it coordinates
people. Fellowship does not require consensus, homogeneity,
reconciliation, or suppressed conflict.

\textbf{Revision conditions.} Institutional expressions must change when
they normalize extraction, silence affected groups, self-certify
correctability, or preserve order by destroying standing and truth.
\end{quote}

\subsection{15.13 A vow of fellowship in the
open}\label{a-vow-of-fellowship-in-the-open}

We commit to build common life without making sameness its price.\\
We commit to tell the truth about who benefits and who bears the cost.\\
We commit to protect dissent without pretending that dissent cannot
err.\\
We commit to submit institutions to evidence and challenge they do not
control.\\
We commit to preserve the material conditions through which agency
becomes real.\\
We commit to remember harms that unity would prefer to forget.\\
We commit to repair what our power has helped damage.\\
And we commit to stories capacious enough to orient a shared future
without becoming sovereign over the people who inhabit them.

\subsection{15.14 Boundaries of the
commitment}\label{boundaries-of-the-commitment}

This account does not license forced harmony, civilizational supremacy,
one-party unity, compulsory spirituality, romanticization of community,
mythic authorization of sacrifice, or the claim that one institution
speaks for humanity.

Fellowship remains fellowship only while persons can remain persons
within it.

\begin{center}\rule{0.5\linewidth}{0.5pt}\end{center}

\section{16. Conclusion: a common world without enforced
sameness}\label{conclusion-a-common-world-without-enforced-sameness}

Social organization is neither an enlarged mind nor a single objective
function. It is the changing pattern through which plural agents
coordinate, contest, depend upon one another, exercise power, and
attempt repair across time.

The framework developed here begins by refusing several shortcuts. Order
is not legitimacy. Compliance is not consent. Trust is not
trustworthiness. Stability is not justice. Average benefit is not
protection of every affected group. A complaint channel is not practical
correctability. A shared story is not legitimate merely because it
unifies.

These refusals make room for a positive account. Collective life becomes
more defensible when institutions preserve independent evidence,
meaningful voice, material capability, protected dissent, accountable
power, practical appeal, and repair. Coordination can be centralized or
distributed depending on information, urgency, scale, and consequence;
no architecture is naturally legitimate. Culture can transmit memory,
care, and belonging while remaining open to reinterpretation and
refusal. Myth can orient long-horizon responsibility while remaining
truth-constrained, plural, and revisable.

Artificial intelligence intensifies every part of the problem. It can
reorganize public knowledge, labor, administration, culture, and
intimacy. It can extend human agency or enclose it. It can help
articulate a common horizon or amplify narrative capture. Its social
meaning will be determined not only by model architecture but by
institutions, ownership, permissions, material infrastructure, public
authorship, and the capacity to correct the systems built around it.

The social task of Coherence Universalism is therefore not to eliminate
difference or discover one final institutional design. It is to help
build a common world in which difference remains speakable, power
remains answerable, truth can still enter, and harm can still call forth
repair.

That is fellowship in the open.

\begin{center}\rule{0.5\linewidth}{0.5pt}\end{center}

\section{Appendix A. Social principles and claim-governance
summary}\label{appendix-a.-social-principles-and-claim-governance-summary}

\subsection{A.1 Domain authority}\label{a.1-domain-authority}

Social Dynamics owns the detailed CU-S, CU-ICF, CU-Inst, CU-EC, and
CU-CL families. Foundations v0.10 records their status and cross-domain
relations without converting every domain principle into a framework
first principle.

The complete wording, evidence grades, participation classes, and review
gates appear in the technical supplement and canonical registry.

\subsection{A.2 Primary social
principles}\label{a.2-primary-social-principles}

\subsubsection{CU-S1 --- Distributed Social
Organization}\label{cu-s1-distributed-social-organization-1}

Social organization consists in the capacity of partially autonomous
agents and institutions to coordinate, contest, correct, and materially
support collective life under disagreement and change. Its assessment is
relative to affected groups, power, protected conditions, scale, and
horizon.

\subsubsection{CU-S2 --- Coordination Through Mutual
Adjustment}\label{cu-s2-coordination-through-mutual-adjustment-1}

Coordination arises through expectations, signals, norms, incentives,
roles, infrastructures, and feedback by which actors adjust to one
another. Coordination may be voluntary, negotiated, habitual, dependent,
or coerced and therefore requires separate legitimacy and distributional
analysis.

\subsubsection{CU-S3 --- Institutional Continuity and
Correctability}\label{cu-s3-institutional-continuity-and-correctability-1}

Institutions encode roles, rules, records, resources, and enforcement
across time. Adequate institutional organization requires not only
continuity but demonstrated capacities to detect error, protect
challenge, provide practically accessible and sufficiently independent
review, reverse decisions, correct records, repair harm, track
recurrence, revise rules, and learn. Nominal channels or
self-certification are insufficient.

\subsubsection{CU-S4 --- Cultural Transmission and
Renewal}\label{cu-s4-cultural-transmission-and-renewal-1}

Culture transmits and transforms meanings, practices, identities,
expectations, and repertoires of coordination across persons and
generations. Cultural continuity is compatible with internal plurality
and requires routes for reinterpretation, refusal, and renewal.

\subsubsection{CU-S5 --- Social Legitimacy Under Plural
Agency}\label{cu-s5-social-legitimacy-under-plural-agency-1}

A social order gains legitimate authority only insofar as those subject
to it retain protected standing and meaningful agency, can understand
and contest materially governing rules, can expose error without
catastrophic retaliation, and have access to effective correction,
repair, or proportionate exit. Stability and popularity are
insufficient.

\subsection{A.3 Canonical civilizational design
extension}\label{a.3-canonical-civilizational-design-extension}

\subsubsection{CU-CL-12 --- Open Myth and Narrative
Pluralism}\label{cu-cl-12-open-myth-and-narrative-pluralism-2}

CU-CL-12 is canonical under the Version 0.10 canonical registry as a
civilizational design principle. Its universal empirical necessity
remains unestablished. Numbering identifies the controlling CU proposal;
it does not complete external validation or authorize one story,
institution, state, or company to govern the meaning of others.

\subsection{A.4 Commitment anchors}\label{a.4-commitment-anchors}

\subsubsection{CU-COM11 --- Fellowship, Justice, and Collective
Repair}\label{cu-com11-fellowship-justice-and-collective-repair-1}

CU-COM11 is the Social Dynamics domain commitment anchor. CU-COM18 names
Heaven on Earth as a revisable social vocation, and CU-COM19 names the
corresponding task of building Omega ecologies. These commitments orient
plural agency, material security, protected dissent, truthful public
records, legitimate authority, love, ecological reciprocity, and
practical correction. They are not descriptive proofs or rules for
selecting one final political regime.

\subsection{A.5 Nineteen application-governance
rules}\label{a.5-nineteen-application-governance-rules}

\begin{enumerate}
\def\labelenumi{\arabic{enumi}.}
\tightlist
\item
  Name the social unit, jurisdiction, and boundary.
\item
  Identify affected groups, including those outside the formal decision
  body.
\item
  Map authority, ownership, agenda-setting power, sanctions, and exit
  costs.
\item
  State the temporal horizon and any discounting of future burdens.
\item
  Declare protected floors separately from weighted benefits.
\item
  Report distributions, not only averages.
\item
  Distinguish trust from trustworthiness and compelled reliance.
\item
  Distinguish compliance, consent, sociological legitimacy, and
  normative legitimacy.
\item
  Preserve independent evidence channels and disclose provenance.
\item
  Protect dissent without treating dissent as automatically true.
\item
  Distinguish coordination capacity from justice or moral value.
\item
  Do not restore a universal social-coherence scalar.
\item
  Declare proxies, missing data, and measurement-induced incentives.
\item
  Test for metric capture, gaming, and classification harm.
\item
  Distinguish association, mechanism, intervention effect, and
  interpretation.
\item
  Do not diagnose populations with individual clinical concepts.
\item
  Give affected communities substantive authority in classifying harms
  and evaluating repair.
\item
  Keep spiritual and civilizational interpretation from substituting for
  empirical or normative argument.
\item
  Do not treat nominal appeals, complaint channels, or self-reported
  review statistics as evidence of correctability without
  practical-access, independence, remedy, recurrence, and affected-party
  evidence.
\end{enumerate}

\section{Appendix B. Major revisions from
v0.09}\label{appendix-b.-major-revisions-from-v0.09}

Version 0.10 retains the integrative project while making the following
changes:

\begin{itemize}
\tightlist
\item
  retires the universal social-coherence functional and general social
  gradient;
\item
  replaces a three-field ontology with a declared multilevel social
  model;
\item
  separates coordination from legitimacy and compliance from consent;
\item
  makes power, distribution, and affected-group scope mandatory;
\item
  replaces universal collapse and attractor claims with conditional
  models;
\item
  distinguishes trust from trustworthiness and dissent from automatic
  truth;
\item
  constrains social-trauma and polarization language against population
  diagnosis;
\item
  treats technology and AI as institutional systems;
\item
  strengthens practical correctability against nominal-channel and
  self-certification failures;
\item
  adds Open Myth as a provisional cultural architecture rather than a
  compulsory worldview;
\item
  preserves the Second Axial Age as an interpretive frame rather than a
  scientific law;
\item
  and places fellowship, justice, and collective repair in the explicit
  commitment register.
\end{itemize}

The complete migration, audit, and patch history is preserved in the
technical archive.

\section{Appendix C. Technical supplement and
reproducibility}\label{appendix-c.-technical-supplement-and-reproducibility}

\textbf{Released evidence package.} The current supplement is
\href{../supplements/07_Social_Dynamics_Technical_Supplement_v0.10.md}{07\_Social\_Dynamics\_Technical\_Supplement\_v0.10.md}.
Its companion evidence directory is
\href{../evidence/07_Social_Dynamics/README.md}{07\_Social\_Dynamics}.
The supplement and evidence directory expose the derivations,
deterministic constructions, scripts, machine-readable results, source
ledger, and audit materials referenced in this manuscript. They do not
elevate a toy model, simulation, or reconstructed check into
target-domain empirical validation.

The accompanying technical supplement contains:

\begin{itemize}
\tightlist
\item
  the formal architecture and non-entailment propositions;
\item
  the deterministic institutional-correction model;
\item
  metric windows and indexing conventions;
\item
  feedback-delay analysis;
\item
  the complete claim and evidence ledger;
\item
  the full provisional social, institutional, economic, and cultural
  registry;
\item
  the Open Myth and mythic--constitutional source architecture;
\item
  the independent audit and bounded repair;
\item
  literature and operationalization ledgers;
\item
  glossary;
\item
  source hashes, code, and deterministic verification results.
\end{itemize}

The constructed model is pedagogical. It does not validate a political
regime or establish the empirical truth of the framework.

\begin{center}\rule{0.5\linewidth}{0.5pt}\end{center}

\section{References}\label{references}

Acemoglu, D., \& Robinson, J. A. (2012). \emph{Why nations fail: The
origins of power, prosperity, and poverty}. Crown.

Anderson, E. (2006). The epistemology of democracy. \emph{Episteme,
3}(1--2), 8--22.

Bail, C. A., Argyle, L. P., Brown, T. W., Bumpus, J. P., Chen, H.,
Hunzaker, M. B. F., Lee, J., Mann, M., Merhout, F., \& Volfovsky, A.
(2018). Exposure to opposing views on social media can increase
political polarization. \emph{Proceedings of the National Academy of
Sciences, 115}(37), 9216--9221.

Bakiner, O. (2021). Truth commission impact on policy, courts, and
society. \emph{Annual Review of Law and Social Science, 17}, 73--91.

Bellah, R. N. (2011). \emph{Religion in human evolution: From the
Paleolithic to the Axial Age}. Belknap Press of Harvard University
Press.

Centola, D. (2018). \emph{How behavior spreads: The science of complex
contagions}. Princeton University Press.

Coleman, J. S. (1990). \emph{Foundations of social theory}. Belknap
Press of Harvard University Press.

Durkheim, É. (1912/1995). \emph{The elementary forms of religious life}
(K. E. Fields, Trans.). Free Press.

Eisenstadt, S. N. (Ed.). (1986). \emph{The origins and diversity of
Axial Age civilizations}. State University of New York Press.

Finkel, E. J., Bail, C. A., Cikara, M., Ditto, P. H., Iyengar, S., Klar,
S., Mason, L., McGrath, M. C., Nyhan, B., Rand, D. G., Skitka, L. J.,
Tucker, J. A., Van Bavel, J. J., Wang, C. S., \& Druckman, J. N. (2020).
Political sectarianism in America. \emph{Science, 370}(6516), 533--536.

Fraser, N. (2008). \emph{Scales of justice: Reimagining political space
in a globalizing world}. Columbia University Press.

Gauthier, G., Hodler, R., Widmer, P., \& Zhuravskaya, E. (2026). The
political effects of X's feed algorithm. \emph{Nature, 652}, 416--423.

Guess, A. M., Malhotra, N., Pan, J., Barberá, P., Allcott, H., Brown,
T., Crespo-Tenorio, A., Dimmery, D., Freelon, D., Gentzkow, M.,
González-Bailón, S., Kennedy, E., Kim, Y. M., Lazer, D., Moehler, D. C.,
Nyhan, B., Rivera, C. V., Settle, J., Thomas, D. R., \ldots{} Tucker, J.
A. (2023a). How do social media feed algorithms affect attitudes and
behavior in an election campaign? \emph{Science, 381}(6656), 398--404.

Guess, A. M., Malhotra, N., Pan, J., Barberá, P., Allcott, H., Brown,
T., Crespo-Tenorio, A., Dimmery, D., Freelon, D., Gentzkow, M.,
González-Bailón, S., Kennedy, E., Kim, Y. M., Lazer, D., Moehler, D. C.,
Nyhan, B., Rivera, C. V., Settle, J., Thomas, D. R., \ldots{} Tucker, J.
A. (2023b). Reshares on social media amplify political news but do not
detectably affect beliefs or opinions. \emph{Science, 381}(6656),
404--408.

Habermas, J. (1996). \emph{Between facts and norms: Contributions to a
discourse theory of law and democracy} (W. Rehg, Trans.). MIT Press.

Hardy, M. D., Thompson, B. D., Krafft, P. M., \& Griffiths, T. L.
(2023). Resampling reduces bias amplification in experimental social
networks. \emph{Nature Human Behaviour, 7}, 2084--2098.

Hartman, R., Blakey, W., Womick, J., Bail, C., Finkel, E. J., Han, H.,
Sarrouf, J., Schroeder, J., Sheeran, P., Van Bavel, J. J., Willer, R.,
\& Gray, K. (2022). Interventions to reduce partisan animosity.
\emph{Nature Human Behaviour, 6}, 1194--1205.

Henrich, J. (2016). \emph{The secret of our success: How culture is
driving human evolution, domesticating our species, and making us
smarter}. Princeton University Press.

Hirschman, A. O. (1970). \emph{Exit, voice, and loyalty: Responses to
decline in firms, organizations, and states}. Harvard University Press.

International Labour Organization. (2024). \emph{World social protection
report 2024--26: Universal social protection for climate action and a
just transition}.

Iyengar, S., Lelkes, Y., Levendusky, M., Malhotra, N., \& Westwood, S.
J. (2019). The origins and consequences of affective polarization in the
United States. \emph{Annual Review of Political Science, 22}, 129--146.

Jaspers, K. (1953). \emph{The origin and goal of history} (M. Bullock,
Trans.). Yale University Press.

Kuran, T. (1995). \emph{Private truths, public lies: The social
consequences of preference falsification}. Harvard University Press.

Landemore, H. (2013). \emph{Democratic reason: Politics, collective
intelligence, and the rule of the many}. Princeton University Press.

Lorenz, J., Rauhut, H., Schweitzer, F., \& Helbing, D. (2011). How
social influence can undermine the wisdom of crowd effect.
\emph{Proceedings of the National Academy of Sciences, 108}(22),
9020--9025.

McCoy, J., Rahman, T., \& Somer, M. (2018). Polarization and the global
crisis of democracy: Common patterns, dynamics, and pernicious
consequences for democratic polities. \emph{American Behavioral
Scientist, 62}(1), 16--42.

Nord, M., Altman, D., Fernandes, T., Good God, A., \& Lindberg, S. I.
(2026). \emph{Democracy report 2026: Unraveling the democratic era?}
V-Dem Institute, University of Gothenburg.

Nussbaum, M. C. (2011). \emph{Creating capabilities: The human
development approach}. Belknap Press of Harvard University Press.

Nyhan, B., Settle, J., Thorson, E., Wojcieszak, M., Barberá, P., Chen,
A. Y., Allcott, H., Brown, T., Crespo-Tenorio, A., Dimmery, D., Freelon,
D., Gentzkow, M., González-Bailón, S., Guess, A. M., Kennedy, E., Kim,
Y. M., Lazer, D., Malhotra, N., Moehler, D. C., \ldots{} Tucker, J. A.
(2023). Like-minded sources on Facebook are prevalent but not
polarizing. \emph{Nature, 620}, 137--144.

OECD. (2024). \emph{OECD survey on drivers of trust in public
institutions---2024 results: Building trust in a complex policy
environment}. OECD Publishing.

Okimoto, T. G., \& Gollwitzer, M. (2025). The social psychology of
justice repair. \emph{Annual Review of Psychology, 76}, 693--716.

Ostrom, E. (1990). \emph{Governing the commons: The evolution of
institutions for collective action}. Cambridge University Press.

Ostrom, E. (2010). Polycentric systems for coping with collective action
and global environmental change. \emph{Global Environmental Change,
20}(4), 550--557.

Page, S. E. (2007). \emph{The difference: How the power of diversity
creates better groups, firms, schools, and societies}. Princeton
University Press.

Putnam, R. D. (2000). \emph{Bowling alone: The collapse and revival of
American community}. Simon \& Schuster.

Rawls, J. (1993). \emph{Political liberalism}. Columbia University
Press.

Rader, G. K. D. (2026). \emph{Coherence Universalism:
Physics---Coherence, Accessibility, and the Ordered Emergence of
Physical Description} (Version 0.10). Heaven≡Earth Foundation.

Rader, G. K. D. (2026). \emph{Coherence Universalism: Metaphysics and
Epistemology---Coherence Under Constraint, Explanatory Order, and
Warrant Across Registers} (Version 0.10). Heaven≡Earth Foundation.

Rader, G. K. D. (2026). \emph{Coherence Universalism:
Foundations---Principles, Domain Gateways, Commitments, and Constraints
on Responsible Application} (Version 0.10). Heaven≡Earth Foundation.

Rader, G. K. D. (2026). \emph{Coherence Universalism: Biology---Life as
Active Self-Maintenance: Viability, Individuation, and Multiscale
Regulation} (Version 0.10). Heaven≡Earth Foundation.

Rader, G. K. D. (2026). \emph{Coherence Universalism:
Consciousness---Perspective, Temporal Depth, Valence, and the
Architecture of Experience} (Version 0.10). Heaven≡Earth Foundation.

Rader, G. K. D. (2026). \emph{Coherence Universalism:
Psychology---Meaning, Regulation, and Repair Across Time} (Version
0.10). Heaven≡Earth Foundation.

Scheffer, M., Bascompte, J., Brock, W. A., Brovkin, V., Carpenter, S.
R., Dakos, V., Held, H., van Nes, E. H., Rietkerk, M., \& Sugihara, G.
(2009). Early-warning signals for critical transitions. \emph{Nature,
461}, 53--59.

Scott, J. C. (1998). \emph{Seeing like a state: How certain schemes to
improve the human condition have failed}. Yale University Press.

Sen, A. (1999). \emph{Development as freedom}. Alfred A. Knopf.

Tainter, J. A. (1988). \emph{The collapse of complex societies}.
Cambridge University Press.

Tyler, T. R. (2006). Psychological perspectives on legitimacy and
legitimation. \emph{Annual Review of Psychology, 57}, 375--400.

Tyler, T. R. (2023). Whither legitimacy? Legal authority in the
twenty-first century. \emph{Annual Review of Law and Social Science,
19}, 1--17.

United Nations Department of Economic and Social Affairs. (2025).
\emph{World social report 2025: A new policy consensus to accelerate
social progress}. United Nations.

Voelkel, J. G., Chu, J., Stagnaro, M. N., Mernyk, J. S., Redekopp, C.,
Pink, S. L., Druckman, J. N., Rand, D. G., \& Willer, R. (2023).
Interventions reducing affective polarization do not necessarily improve
anti-democratic attitudes. \emph{Nature Human Behaviour, 7}, 55--64.

WHO Commission on Social Connection. (2025). \emph{From loneliness to
social connection: Charting a path to healthier societies}. World Health
Organization.

Woolley, A. W., Chabris, C. F., Pentland, A., Hashmi, N., \& Malone, T.
W. (2010). Evidence for a collective intelligence factor in the
performance of human groups. \emph{Science, 330}(6004), 686--688.

World Bank. (2024). \emph{A framework for understanding and measuring
resilience}. World Bank Social Sustainability and Inclusion Global
Practice.

Young, I. M. (2000). \emph{Inclusion and democracy}. Oxford University
Press.

Gokhale, C. S., Bulbulia, J., \& Frean, M. (2022). Collective narratives
catalyse cooperation. \emph{Humanities and Social Sciences
Communications, 9}, 85.

Rader, G. K. D. (2026). \emph{The Book of Eight Gates}. Heaven≡Earth
Press.

\end{document}
